\section{Numbered statement register}\label{sec:K}
This register lists every current result row, including unresolved parent tasks.
The identifiers are stable across the report, calculations and check-event table.
The full typed statements, diplomatic French quotations and proofs occur at the
locations given below. A proved special case does not change an unresolved parent
task to proved. Historical attempts are retained with their dated verification state.
\begingroup\small\raggedright
\par\medskip\noindent\textbf{K.1. A1 --- Undecided.}\par
Picard isomorphism for every smooth hyperplane in dimension at least four.
\par\emph{Statement/proof location:} Part A, A1r and A4.
\emph{Manuscript reference:} 57:97.
\par Reduced to the dimension-four case; finite multiplicative obstruction remains.
\par\medskip\noindent\textbf{K.2. A2 --- Undecided.}\par
Injectivity and the least universal dimension.
\par\emph{Statement/proof location:} Part A, A2b and A4b.
\emph{Manuscript reference:} 57:97.
\par Universal point injectivity and its least dimension remain undecided. The kernel in dimension at least three is a finite p-group; for each nonzero order-p class its bad smooth hyperplanes form the finite determinantal locus A2D. A2V forces a Cartier-exact leading term and Frobenius kernel on negative H1 of the hyperplane, with vanishing order divisible by p and a nonnormal mu\_p torsor. A2W forces rank(d:H0Omega1->H0Omega2)>=n-1. A2Y additionally requires noninjective Frobenius on H1(O\_Y) and no smooth W2 lifting of Y. In dimension3 the target is of general type by applying A3s to L\textasciicircum{}a. No global realization or universal exclusion of this remaining gate is proved.
\par\medskip\noindent\textbf{K.3. A3 --- Proved.}\par
There exists a smooth projective surface with a very ample L and nonzero H1(L\textasciicircum{}-1); the least requested dimension is two.
\par\emph{Statement/proof location:} sections/A.tex, a:very-ample-example.
\emph{Manuscript reference:} 57:99-100.
\par Explicit characteristic-two Raynaud surface followed by 556 simultaneous ninth-root covers; h0(L)=562 and h1(L\textasciicircum{}-1)=1. Complete GF(2) basis and finite transition certificate included; three targeted independent audits passed.
\par\medskip\noindent\textbf{K.4. A4 --- Proved.}\par
Exact finite multiplicative obstruction description of kernel and cokernel.
\par\emph{Statement/proof location:} Part A, Statement a:finite.
\emph{Manuscript reference:} 57:97.
\par Uses the actual multiplicative tower; not just dimensions or a direct sum of additive layers.
\par\medskip\noindent\textbf{K.5. B1 --- Refuted.}\par
Seshadri without Pic as formulated in prompt.
\par\emph{Statement/proof location:} work/bce/section\_B.tex; section.tex.
\emph{Manuscript reference:} 57:110.
\par Split projective quadric xy=z(z-tw), D=(x,z), all characteristics; original manuscript retains essential properness of Pic functor
\par\medskip\noindent\textbf{K.6. B2 --- Refuted as a fibre-type-only criterion; fixed-family iff proved.}\par
Exact conditions in Question 2 and necessity of R2.
\par\emph{Statement/proof location:} sections/B.tex, B2 exact obstruction and paired projective families.
\par Identical regular base and abstract geometric fibre types give opposite extension properties; no criterion from those data alone exists. The exact deformation-sensitive Cl/Pic kernel criterion is proved for regular base and normal integral geometric fibres. R2 is not necessary and requires the original lci hypothesis for sufficiency.
\par\medskip\noindent\textbf{K.7. C1 --- Proved.}\par
Projectivity descends from completion.
\par\emph{Statement/proof location:} work/bce/section\_C.tex; Artin1969 Thm3.5 pp31-32.
\emph{Manuscript reference:} 57:227-231;272.
\par Artin approximation of an ample bundle on the closed fibre; relative ampleness over a local base
\par\medskip\noindent\textbf{K.8. C2 --- Proved.}\par
Picard comparison when R1 has finite length.
\par\emph{Statement/proof location:} sections/C.tex, Theorem C2.G.
\par Full excellent henselian comparison proved by principal finite determination, blow-up and Stein factorization, conductor descent, and approximation; independent proof audit passed.
\par\medskip\noindent\textbf{K.9. C3 --- Proved.}\par
Punctured henselian Pic equals completed Pic.
\par\emph{Statement/proof location:} work/bce/section\_C.tex; Elkik1973 III.1 Thm3 p569 and Pic cor p573.
\emph{Manuscript reference:} 57:194;227-231.
\par Any noetherian henselian pair; no isolated-singularity or characteristic restriction
\par\medskip\noindent\textbf{K.10. C4 --- Proved in original cone scope; arbitrary-local extrapolation refuted.}\par
Local Neron-Severi comparison modulo p groups.
\par\emph{Statement/proof location:} sections/C.tex, cone comparison and general-local counterexample.
\par The source fixes the local cone at its vertex and then compares punctured Picard groups. That statement is proved modulo bounded p-groups. A characteristic-three factorial local hypersurface with henselian/completed punctured Pic Z refutes the arbitrary-local extrapolation; different intrinsic etale-local NS definitions are stated separately.
\par\medskip\noindent\textbf{K.11. D1-positive --- Proved.}\par
Positive quotient Hodge asymmetry.
\par\emph{Statement/proof location:} work/dgj/section\_D.tex.
\emph{Manuscript reference:} 57:287.
\par Characteristic two shearing involutions with F and V; dimensions (2,3),(3,2), specializing to (3,3).
\par\medskip\noindent\textbf{K.12. D1-negative --- Proved.}\par
Negative quotient Hodge asymmetry.
\par\emph{Statement/proof location:} work/dgj/section\_D.tex.
\emph{Manuscript reference:} 57:287.
\par Characteristic two shearing involutions with F and V; dimensions (2,3),(3,2), specializing to (3,3).
\par\medskip\noindent\textbf{K.13. D1-family --- Proved.}\par
Variation of quotient Hodge asymmetry.
\par\emph{Statement/proof location:} work/dgj/section\_D.tex.
\emph{Manuscript reference:} 57:287.
\par Characteristic two shearing involutions with F and V; dimensions (2,3),(3,2), specializing to (3,3).
\par\medskip\noindent\textbf{K.14. D2-flatness --- Proved.}\par
Nonflat fixed groups and Picard torsion.
\par\emph{Statement/proof location:} work/dgj/section\_D.tex.
\emph{Manuscript reference:} 57:288.
\par Fixed lengths 8 and 16; algebra R[z]/(z\textasciicircum{}4,t z\textasciicircum{}2).
\par\medskip\noindent\textbf{K.15. D2-specialization --- Proved.}\par
Picard variety and specialization.
\par\emph{Statement/proof location:} work/dgj/section\_D.tex.
\emph{Manuscript reference:} 57:288.
\par Fibrewise image of (F\textasciicircum{}2,[2]) has kernel degree 2 then 4.
\par\medskip\noindent\textbf{K.16. D2-Igusa --- Proved.}\par
Modular family of the Igusa surface.
\par\emph{Statement/proof location:} work/dgj/section\_D.tex D4o and D4s; section\_D\_supersingular.tex.
\emph{Manuscript reference:} 57:311.
\par Both algebraized miniversal hulls computed: ordinary W[[t,s,u,v]]/(s\ensuremath{{}^2}+2s); supersingular W[[t,s,u,v]]/(s\ensuremath{{}^4}−t*s\ensuremath{{}^2}−2s). Canonical fixed elliptic subgroup and all-Artinian Pic-tau flatness proved; root/fhi independent audits pass.
\par\medskip\noindent\textbf{K.17. E1 --- Refuted as stated; intended valuative version proved.}\par
Properness of Pic0 over a trait.
\par\emph{Statement/proof location:} sections/E.tex, E1 literal counterexample and intended valuative theorem.
\par A normal total-space elliptic degeneration refutes the literal assertion. The version with geometrically normal fibres is proved over a trait. Arbitrary-base representability questions are distinct E4 rows.
\par\medskip\noindent\textbf{K.18. E2 --- Proved.}\par
Corr is separated under proper flat universal H0=O.
\par\emph{Statement/proof location:} work/bce/section\_E.tex; Raynaud1970 3.2(c),5.2;Murre1966 Thm4.
\emph{Manuscript reference:} 57:22;32;313.
\par Trait normalization into étale closure E of identity; connected fibres force map zero; Murre yields unramified separated scheme
\par\medskip\noindent\textbf{K.19. E3 --- Proved.}\par
Problem A universal flattening monomorphism of finite presentation.
\par\emph{Statement/proof location:} work/bce/section\_E.tex; Stacks05PS Lem38.27.6(3);0CWQ Lem77.11.8(3).
\emph{Manuscript reference:} 57:73-79.
\par Proper finitely presented X and finitely presented F; all base changes
\par\medskip\noindent\textbf{K.20. E4 --- Undecided.}\par
Remaining structure questions on 57:308-315.
\par\emph{Statement/proof location:} sections/E.tex, individual questions 57:308-315.
\par E4.T proves Pic-tau closedness without projectivity for proper flat finitely presented geometrically integral families over arbitrary scheme bases. E4.S settles the fibrewise smoothness/defect claims; E4.Curve proves full Pic schematicity for smooth proper geometrically connected curves. Remaining mathematical gate: full Pic schematicity for geometrically integral nonnormal families and normal higher-dimensional families over arbitrary bases. The scan310 finite-component iff retains an explicitly recorded reading/scope uncertainty; multiplication nonopenness is already refuted by E4.H.
\par\medskip\noindent\textbf{K.21. F --- Undecided.}\par
Finite generation of the full local etale fundamental group.
\par\emph{Statement/proof location:} work/fhi/F\_proof.md section F4.
\emph{Manuscript reference:} 58:19.
\par For complete normal local surfaces with algebraically closed residue field, the unrestricted conjecture remains. In equal characteristic F8/F9b settle every geometric quotient with normal pro-p kernel over the full tame quotient. F10 proves finite abelian inertia using Krishna–Majumder2026v2; F10c gives all-prime rank-linear H1 bounds for abelian monodromy, and F11 proves uniform finite generation of metabelian quotients. F12 gives the characteristic-p bound for metabelian coefficients; F12b settles pro-p closed second derived subgroup. F13 proves a fixed log-conductor bound after one fixed cover. F13c/d remove unramified character order and reduce the remaining coefficient-prime-not-p estimate to a degree-linear tame-generator bound for cyclic covers of uniformly bounded p-power degree (F13.6). F13e bounds refined-Swan zero length, but not cleanup complexity. F13.5/F13.6, further solvable and nonsolvable wild quotients, and mixed characteristic remain unresolved. The actual wild kernel is not assumed pro-p.
\par\medskip\noindent\textbf{K.22. G1 --- Proved.}\par
Universal Pic-delta section exists iff at most one elementary divisor is odd.
\par\emph{Statement/proof location:} work/dgj/section\_G.tex.
\emph{Manuscript reference:} 59:10-14.
\par Over Z[1/2], construction uses rank of induced pairing on A[2].
\par\medskip\noindent\textbf{K.23. G2 --- Proved.}\par
Determinant constructions agree with canonical sections.
\par\emph{Statement/proof location:} work/dgj/section\_G.tex.
\emph{Manuscript reference:} 59:12-14.
\par Includes square and even-Euler double transform comparisons; corrects ppcm and even-chi divisibility claim.
\par\medskip\noindent\textbf{K.24. G3-Albanese --- Proved.}\par
Sign of the Albanese torsor.
\par\emph{Statement/proof location:} work/dgj/section\_G.tex.
\emph{Manuscript reference:} 59:23-28.
\par Minus sign visually verified; canonical norm on PPAV theta torsor.
\par\medskip\noindent\textbf{K.25. G3-principal --- Proved.}\par
Canonical class of principally polarized abelian schemes.
\par\emph{Statement/proof location:} work/dgj/section\_G.tex.
\emph{Manuscript reference:} 59:23-28.
\par Minus sign visually verified; canonical norm on PPAV theta torsor.
\par\medskip\noindent\textbf{K.26. H-Ga --- Proved.}\par
Higher Ext vanishing for the additive group.
\par\emph{Statement/proof location:} work/fhi/H.tex; H\_proof.md.
\emph{Manuscript reference:} 43:80.
\par Corrects assertion that Ga remains difficult
\par\medskip\noindent\textbf{K.27. H-all --- Proved.}\par
Higher Ext vanishing for every algebraic commutative group.
\par\emph{Statement/proof location:} work/fhi/H.tex; H\_proof.md.
\emph{Manuscript reference:} 43:24-35;43:80.
\par All proalgebraic targets; all p
\par\medskip\noindent\textbf{K.28. I-2d --- Proved.}\par
Residual complex endomorphisms in dimension two, degree two.
\par\emph{Statement/proof location:} work/fhi/I.tex; I\_proof.md.
\emph{Manuscript reference:} 43:65-68;SA19-H3.
\par Precise finite-topology universal interpretation; mixed arrow topology
\par\medskip\noindent\textbf{K.29. I-all --- Undecided.}\par
Completed algebra presentation in every dimension.
\par\emph{Statement/proof location:} work/fhi/I\_universal\_linear.tex; I\_proof.md.
\emph{Manuscript reference:} 43:65-68;SA19-H3.
\par All dimensions are proved for separated complete graded algebras with a basis of open left ideals; all surfaces are proved for arbitrary jointly continuous topology. The original43:66 specifies the broader category without the open-left-ideal hypothesis. I7b now proves open dense differential-image quotients for whole source columns over arbitrary regular local rings. Remaining: construction and continuity of the entire adjacent-arrow family and its endpoint assembly in arbitrary complete targets, in dimension at least3.
\par\medskip\noindent\textbf{K.30. J1-Th1 --- Proved for common incidence or regular-test kernels; manuscript identification unresolved.}\par
Set-theoretic kernel equality for manuscript Albanese equivalence.
\par\emph{Statement/proof location:} work/dgj/section\_J.tex J1i; section\_J\_incidence.tex.
\emph{Manuscript reference:} 44:2.
\par Common incidence kernels agree scheme-theoretically with the represented family map. Common all-regular kernels agree on points. Undefined manuscript N and en substance/ensemblistement variants remain explicit.
\par\medskip\noindent\textbf{K.31. J1-Th2 --- Known and proved for incidence; all-regular generality unresolved.}\par
Bound for all manuscript quotients.
\par\emph{Statement/proof location:} work/dgj/section\_J.tex J1i; section\_J\_incidence.tex.
\emph{Manuscript reference:} 44:2.
\par Every incidence family quotient has dimension at most b(2i-1)/2 in all codimensions and characteristics. For all regular maps, general uniform boundedness is equivalent to universal existence, with codimension3/fourfold gate.
\par\medskip\noindent\textbf{K.32. J1-Th3 --- Incidence existence and isogeny duality known; modern integral duality refuted; general regular existence unresolved.}\par
General representability and complementary integral duality.
\par\emph{Statement/proof location:} work/dgj/section\_J.tex J1i and J2a.
\emph{Manuscript reference:} 44:2;44:8.
\par Saito constructs the incidence quotient in all codimensions with specified isogeny normalization. The modern Ab integral assertion is false. Kahn2021 Conjecture1 gives a concrete remaining regular-representability test.
\par\medskip\noindent\textbf{K.33. J1-Th4 --- Refuted.}\par
Raw middle intersection autoduality is positive up to one sign.
\par\emph{Statement/proof location:} work/dgj/section\_J.tex.
\emph{Manuscript reference:} 44:2;44:8-9.
\par Primitive-qualified question is not refuted.
\par\medskip\noindent\textbf{K.34. J1-Th5 --- Refuted.}\par
Higher codimension height/intersection comparison as revised.
\par\emph{Statement/proof location:} work/dgj/section\_J.tex J5; section\_J5\_quadric.tex.
\emph{Manuscript reference:} 44:4.
\par Over Fbar3: Q4, pencil x2=t*x5, all fibres geometrically integral. Both restricted cycles literally zero; Z.Zprime=2 from two transverse points. i=j=m=2 fixes the sign. Excess intersections with the pencil base are permitted by the literal statement; stronger admissibility conditions remain distinct.
\par\medskip\noindent\textbf{K.35. J1-Th6 --- Refuted as stated in the middle clause; full stable and high clauses unresolved.}\par
Lefschetz for all algebraic representatives.
\par\emph{Statement/proof location:} work/dgj/section\_J.tex J6C/J6L/J6P/J6Z; section\_J6\_cubic.tex.
\emph{Manuscript reference:} 44:4.
\par J6C: over Fbar5 a pencil with every fibre geometrically integral has zero ambient Ab2 and Pinc2 but a constant elliptic factor in the generic member, detected by composition[-2]. No smooth-axis hypothesis occurs in44:4. J6P proves the smooth-axis codimension-two correction; J6L and J6Z give conditional/general and unconditional top-target cases. Full stable/high, positive-characteristic finite-kernel and integral-isomorphism readings remain gates.
\par\medskip\noindent\textbf{K.36. J1-Th7 --- Refuted.}\par
Numerical equals tau equivalence for all cycles.
\par\emph{Statement/proof location:} work/dgj/section\_J.tex.
\emph{Manuscript reference:} 44:6.
\par Explicit Fermat quartic Ceresa cycle over Qbar; Bloch1984 pp94-95 and Theorem4.1 p104 visually verified. Extends to all intermediate codimensions by products.
\par\medskip\noindent\textbf{K.37. J2-isomorphism --- Refuted.}\par
Threefold algebraic representative canonical self-duality.
\par\emph{Statement/proof location:} work/dgj/section\_J.tex.
\emph{Manuscript reference:} 44:2;44:8-9.
\par Generic monic septic Jacobian over the algebraic closure of F5(a0,...,a6)\textasciicircum{}perf; intersection forces a degree8 K-isogeny, ruled out by irreducible S7 action on A[2]. No integral torsion standard assumption used.
\par\medskip\noindent\textbf{K.38. J2-positivity --- Refuted.}\par
Threefold algebraic representative polarization.
\par\emph{Statement/proof location:} work/dgj/section\_J.tex.
\emph{Manuscript reference:} 44:2;44:8-9.
\par Explicit E/F3 times P1 times P1; matrix [[0,lambda],[lambda,0]].
\par\medskip\noindent\textbf{K.39. K --- Current register complete; final audit pending.}\par
Complete numbered statement of results.
\par\emph{Statement/proof location:} sections/K.tex, complete current numbered register.
\par Every current A-J and lead row is numbered with status and proof location. Final completeness requires resolution or the user-defined stopping review for every remaining research gate.
\par\medskip\noindent\textbf{K.40. SA14-H1 --- Proved.}\par
Generic Pic-tau restriction kernel and existence of its nonzero infinitesimal part.
\par\emph{Statement/proof location:} sections/A.tex, a:generic-unipotent and a:very-ample-example.
\emph{Manuscript reference:} 57:99-100.
\par For normal projective X the geometric-generic Pic-tau restriction kernel is finite connected unipotent; its tangent space is H1(L\textasciicircum{}-1). A3 gives an explicit nonzero example in dimension two.
\par\medskip\noindent\textbf{K.41. SA14-H2 --- Refuted as prompt formulated.}\par
Seshadri without Pic.
\par\emph{Statement/proof location:} work/bce/additional\_leads.tex; section.tex.
\emph{Manuscript reference:} 57:110.
\par Same as B1; omitted essential properness in source must be restored
\par\medskip\noindent\textbf{K.42. SA14-H3 --- Refuted as recalled; target undecided.}\par
Universal surjectivity without cohomological vanishing and recalled local analogue.
\par\emph{Statement/proof location:} sections/A.tex, A1r and tested approaches.
\emph{Manuscript reference:} 57:97.
\par Bhatt-de Jong local theorem is injectivity up to p-power torsion, not the recalled isomorphism; no global solution follows.
\par\medskip\noindent\textbf{K.43. SA14-H4 --- Undecided.}\par
Boxed universal injectivity in dimension at least four.
\par\emph{Statement/proof location:} sections/A.tex, A2b and A4b.
\emph{Manuscript reference:} 57:97.
\par Universal point injectivity and its least dimension remain undecided. The kernel in dimension at least three is a finite p-group; for each nonzero order-p class its bad smooth hyperplanes form the finite determinantal locus A2D. A2V forces a Cartier-exact leading term and Frobenius kernel on negative H1 of the hyperplane, with vanishing order divisible by p and a nonnormal mu\_p torsor. A2W forces rank(d:H0Omega1->H0Omega2)>=n-1. A2Y additionally requires noninjective Frobenius on H1(O\_Y) and no smooth W2 lifting of Y. In dimension3 the target is of general type by applying A3s to L\textasciicircum{}a. No global realization or universal exclusion of this remaining gate is proved.
\par\medskip\noindent\textbf{K.44. SA14-H5 --- Proved.}\par
All henselian claims.
\par\emph{Statement/proof location:} work/bce/additional\_leads.tex; section\_C.tex;Artin-Popescu;Elkik1973.
\emph{Manuscript reference:} 57:227-231;272.
\par C1 projectivity, full C2.G Picard completion for arbitrary excellent henselian local base and proper X with finite-length H1, and C3 punctured Pic by Elkik.
\par\medskip\noindent\textbf{K.45. SA14-H6 --- Proved with Pic algebraic-space hypothesis.}\par
Pic/Pic\ensuremath{\tau} exists as an algebraic space if tau flat fp; stronger closed version is schematic.
\par\emph{Statement/proof location:} work/bce/additional\_leads.tex; Stacks06PH;0418.
\emph{Manuscript reference:} 57:277.
\par Flat fp translation quotient is an algebraic space without closedness (prompt summary). HUNCHES additionally assumes tau closed: since tau is open, the quotient is separated unramified and a scheme. Both formulations preserved.
\par\medskip\noindent\textbf{K.46. SA14-H7 --- Unsettled as composite.}\par
Normal trait properness and scheme representability.
\par\emph{Statement/proof location:} work/bce/additional\_leads.tex; section.tex.
\emph{Manuscript reference:} 57:314-315.
\par Literal total-normal assertion false; intended connected relative-normal trait valuation proved; arbitrary-base geomnormal Pic schematicity not yet settled
\par\medskip\noindent\textbf{K.47. SA14-H8 --- Proved with complete dimension correction.}\par
Exceptional divisor lattice injects into NS exceptional fibre.
\par\emph{Statement/proof location:} work/bce/additional\_leads.tex; KollarFamilies Lem11.60 p436.
\emph{Manuscript reference:} 57:8.
\par Kernel exactly free lattice on one-dimensional branches; zero if every component has dimension>=2; use ±D nef and normalization
\par\medskip\noindent\textbf{K.48. SA14-H9 --- Proved.}\par
Corr always separated.
\par\emph{Statement/proof location:} work/bce/section\_E.tex; Raynaud1970 3.2(c),5.2;Murre1966 Thm4.
\emph{Manuscript reference:} 57:22;32;313.
\par Trait normalization into étale closure E of identity; connected fibres force map zero; Murre yields unramified separated scheme
\par\medskip\noindent\textbf{K.49. SA14-H10 --- Proved.}\par
SA14 lead H10.
\par\emph{Statement/proof location:} work/bce/additional\_leads.tex; Deligne1968 Thm5.5(i) p123;SGA3 VI\_B Prop1.6 p327.
\emph{Manuscript reference:} 57:64.
\par Deligne Thm5.5 gives conormal=(R1f*O)\textasciicircum{}dual; SGA3 criterion yields smooth Pic0; proper neutral locus yields abelian scheme
\par\medskip\noindent\textbf{K.50. SA14-H11 --- Proved.}\par
SA14 lead H11.
\par\emph{Statement/proof location:} work/bce/additional\_leads.tex; SGA3 XVI4.2 pp521-522;VI\_B1.4.2 pp325-326;SGA6 XII1.1-1.2.
\emph{Manuscript reference:} 57:74;78.
\par Kernel quotient and fibrewise-finite monomorphism argument, full proof Theorem H11.G.
\par\medskip\noindent\textbf{K.51. SA14-H12 --- Proved in original cone scope; refuted as arbitrary-local extrapolation.}\par
Local NS invariance modulo p-groups.
\par\emph{Statement/proof location:} work/bce/additional\_leads.tex; section.tex.
\emph{Manuscript reference:} 57:193-194.
\par 57:184–192 explicitly fixes cone vertex S;193 compares actual punctured Pic groups modulo identity components,194 first introduces henselization. Cone comparison is modulo bounded p-power. Arbitrary punctured-NS comparison fails 0\ensuremath{\to}Z\ensuremath{\to}Z in char3; a strict-henselian definition names a different invariant.
\par\medskip\noindent\textbf{K.52. SA14-H13 --- Proved/refuted components.}\par
R2 refinement and necessity.
\par\emph{Statement/proof location:} work/bce/additional\_leads.tex; section.tex.
\emph{Manuscript reference:} 57:111.
\par Sufficient with lci regular base; unnecessary; without lci not sufficient
\par\medskip\noindent\textbf{K.53. SA19-H1 --- Known; Undetermined chronology.}\par
Component-group pairing identified with SGA7 IX1.3 and Suzuki perfect-residue theorem; priority date not established.
\par\emph{Statement/proof location:} work/fhi/priority.md; priority.tex.
\emph{Manuscript reference:} 43:117; SA19-H1.
\par Bibliographic dates on 87 and144 and administrative list118 do not date scientific leaf117
\par\medskip\noindent\textbf{K.54. SA19-H2 --- Known identification; Undetermined chronology.}\par
Three-variable Ext matches biextension formalism; claimed pre-SGA7 priority not established.
\par\emph{Statement/proof location:} work/fhi/priority.md; priority.tex.
\emph{Manuscript reference:} 43:16;43:21; SA19-H2.
\par SGA7VII p133 attributes introduction to Mumford; precise archival dating remains gate
\par\medskip\noindent\textbf{K.55. SA19-H3 --- Undecided.}\par
SA19 lead H3.
\par\emph{Statement/proof location:} See the item discussion; no completed proof is claimed..
\par All dimensions are proved for separated complete graded algebras with a basis of open left ideals; all surfaces are proved for arbitrary jointly continuous topology. The original43:66 specifies the broader category without the open-left-ideal hypothesis. I7b now proves open dense differential-image quotients for whole source columns over arbitrary regular local rings. Remaining: construction and continuity of the entire adjacent-arrow family and its endpoint assembly in arbitrary complete targets, in dimension at least3.
\par\medskip\noindent\textbf{K.56. SA19-H4 --- Refuted.}\par
SA19 lead H4.
\par\emph{Statement/proof location:} See the item discussion; no completed proof is claimed..
\par The actual ordinary local Jacobian of k[s,t]\_(s,t) has an explicit nonsplit Z/p extension and is not projective; higher Ext vanishing for the curve pro-Jacobian is separately proved in Serre category and on the perfect etale site.
\par\medskip\noindent\textbf{K.57. SA19-H5 --- Refuted as stated; corrected programme unresolved.}\par
The sweeping open-status classification of the seven theorems.
\par\emph{Statement/proof location:} work/dgj/section\_J.tex.
\emph{Manuscript reference:} 44:2-9.
\par Universal integral intersection autoduality, raw positivity, literal higher-codimension height formula, middle fixed-part clause without smooth pencil axis, and numerical=tau are refuted. Primitive normalization, admissible height corrections, and manuscript-equivalence identification remain separate. Incidence existence/bound are known in every codimension. General regular representability and full stable/high Lefschetz assertions remain unresolved.
\par\medskip\noindent\textbf{K.58. SA19-H6 --- Complete scoped audit; inferred pi1 injectivity refuted.}\par
No readable rational-point-to-section injectivity proof in the supplied cote43 corpus; character-variance correction.
\par\emph{Statement/proof location:} work/dgj/section\_J.tex H6a-H6b; audit\_SA19H6.md; outputs/c43\_audit.tsv.
\emph{Manuscript reference:} 43:1-208 supplied readable transcriptions; exact positive locators43:27-35.
\par 184 complete files read by DGJ and11 by FHI. No sought proof found;2620 uncertainty macros,48 illegibility macros and13 absent scan labels limit the conclusion. Ext1 injectivity at43:30 detects surjectivity of the fundamental-group map, not injectivity. Original scans and historical nonexistence not established.
\par\medskip\noindent\textbf{K.59. A4b --- Proved.}\par
Finiteness and explicit p-power exponent for the restriction kernel and cokernel.
\par\emph{Statement/proof location:} Part A, Corollary a:pgroups.
\emph{Manuscript reference:} 57:97.
\par Both groups for n>=4; kernel for n=3. Exponent p\textasciicircum{}e for p\textasciicircum{}e>=N+1.
\par\medskip\noindent\textbf{K.60. A4c --- Known.}\par
Picard restriction under the stated negative-cohomology vanishings.
\par\emph{Statement/proof location:} Part A, Corollary a:vanishing.
\emph{Manuscript reference:} 57:97.
\par Hypotheses verified for smooth X when the stated vanishings hold; not supplied in A1.
\par\medskip\noindent\textbf{K.61. A3r --- Proved.}\par
Any very ample H1-negative counterexample yields a smooth surface counterexample.
\par\emph{Statement/proof location:} Part A, Proposition a:surface.
\emph{Manuscript reference:} 57:100.
\par Strict reduction of the existence gate; not an existence theorem.
\par\medskip\noindent\textbf{K.62. A1r --- Proved.}\par
The universal dimension>=4 Picard isomorphism reduces to dimension four.
\par\emph{Statement/proof location:} Part A, Proposition a:fourfold.
\emph{Manuscript reference:} 57:97.
\par Auxiliary high-degree cut preserves the original hyperplane embedding on the cut.
\par\medskip\noindent\textbf{K.63. A2b --- Proved.}\par
A smooth Segre surface hyperplane has restriction kernel Z(1,-1).
\par\emph{Statement/proof location:} Part A, Proposition a:segre.
\emph{Manuscript reference:} 57:97.
\par Valid in every characteristic; proves only lower bound three for a possible universal threshold.
\par\medskip\noindent\textbf{K.64. R3 --- Proved.}\par
Tangent of the Picard restriction kernel equals h H1(X,L\textasciicircum{}-1).
\par\emph{Statement/proof location:} Part A, Proposition a:tangent.
\emph{Manuscript reference:} 57:99-100.
\par The equivalence K=0 iff tangent=0 requires K finite connected; generic connectedness audited separately.
\par\medskip\noindent\textbf{K.65. R4 --- Proved.}\par
Restriction kernel/cokernel are bounded finite p-groups in the dimensions stated.
\par\emph{Statement/proof location:} Part A, Corollary a:pgroups.
\emph{Manuscript reference:} 57:97.
\par The phrase built from layers requires obstruction maps and extension data, not direct sums.
\par\medskip\noindent\textbf{K.66. B2-necessity --- Refuted.}\par
Is R2 necessary?.
\par\emph{Statement/proof location:} work/bce/section\_B.tex; section.tex.
\emph{Manuscript reference:} 57:111.
\par Projective regular smoothing xy-z\ensuremath{{}^2}-tw\ensuremath{{}^2}=0 has non-R2 special quadric cone but locally factorial total space; compare xy-z\ensuremath{{}^2}-t\ensuremath{{}^2}w\ensuremath{{}^2}=0 over same trait
\par\medskip\noindent\textbf{K.67. B2-sufficiency --- Proved with original lci hypothesis.}\par
R2 fibres over regular base suffice.
\par\emph{Statement/proof location:} work/bce/section\_B.tex; SGA2 XI 3.13(ii).
\emph{Manuscript reference:} 57:111.
\par Parafactoriality in local dimension >=4; without lci explicit Segre-cone projective family refutes it
\par\medskip\noindent\textbf{K.68. B2-exact --- Proved impossibility of a fibre-type-only criterion; fixed-family iff proved.}\par
Determine exact condition on fibres and S.
\par\emph{Statement/proof location:} work/bce/section\_B.tex; section.tex.
\emph{Manuscript reference:} 57:111.
\par Proper lci families X1 and X2 have identical base and geometric fibre types but opposite extension properties. Thus these data cannot determine the answer. For regular S and normal integral geometric fibres the exact deformation-sensitive obstruction is kernel Cl/Pic\ensuremath{\to}Cl\_generic/Pic\_generic.
\par\medskip\noindent\textbf{K.69. C2-char0 --- Proved in characteristic zero.}\par
Pic(X)=Pic(completion) when R1 has finite length.
\par\emph{Statement/proof location:} work/bce/section\_C.tex; Artin1969 Thm3.5;PerlingSchroer Prop1.2 proof.
\emph{Manuscript reference:} 57:227-231;272.
\par Eventual H1(m\textasciicircum{}n O)=0; formal log; finite determination plus Artin density
\par\medskip\noindent\textbf{K.70. C2-mixed --- Proved in all characteristics for every excellent henselian local base.}\par
Same over excellent DVR bases.
\par\emph{Statement/proof location:} work/bce/section\_C.tex; section\_C.tex Theorem C2.G;Stacks02OB(3),02OS,03H0;PerlingSchroer Prop1.2 E=O;Artin-Popescu.
\emph{Manuscript reference:} 57:227-231;272.
\par Full Theorem C2.G: blow up maximal ideal, Stein factor gives Cartier ideal with R1 supported exceptionally; principal-ideal finite determination, auxiliary Stein H1 and H2 finite length, conductor-unit descent via uniform formal functions, Artin approximation. No flatness assumption on X.
\par\medskip\noindent\textbf{K.71. C4-cones --- Proved in original cone context.}\par
Local cone comparison modulo p-groups.
\par\emph{Statement/proof location:} work/bce/section\_C.tex; section.tex.
\emph{Manuscript reference:} 57:193-194.
\par Split formal neighbourhood; finite filtration by additive H1(O\_X(n)) killed by p; Serre vanishing bounds exponent
\par\medskip\noindent\textbf{K.72. C4-general --- Scope-dependent refutation.}\par
General local NS comparison modulo p-groups.
\par\emph{Statement/proof location:} work/bce/section\_C.tex; section.tex.
\emph{Manuscript reference:} 57:193-194.
\par Char3 factorial local hypersurface xy=z\ensuremath{{}^2}+t\ensuremath{{}^2}(1+t) has completed/henselian punctured Pic Z; intrinsic étale-local NS conventions must be distinguished
\par\medskip\noindent\textbf{K.73. E1-literal --- Refuted.}\par
Normal total space implies extension of algebraically trivial generic bundle.
\par\emph{Statement/proof location:} work/bce/section\_E.tex; section.tex.
\emph{Manuscript reference:} 57:314-315.
\par Normal projective cubic Y\ensuremath{{}^2}Z=X\ensuremath{{}^3}+X\ensuremath{{}^2}Z+t\ensuremath{{}^2}Z\ensuremath{{}^3}; P-O generic degree zero does not extend
\par\medskip\noindent\textbf{K.74. E1-relative --- Proved in intended valuative sense.}\par
Proper flat geometrically normal connected family over trait.
\par\emph{Statement/proof location:} work/bce/section\_E.tex; Raynaud1970 Thm3.3.1;FGA VI Cor2.3.
\emph{Manuscript reference:} 57:314-315.
\par Separated Pic algebraic space is scheme over trait; proper neutral subset yields extension with special class in Pic0
\par\medskip\noindent\textbf{K.75. E1-functor --- Refuted in general.}\par
Pic0 subfunctor is a proper algebraic space.
\par\emph{Statement/proof location:} work/bce/section\_E.tex; LaurentSchroer2024 Prop2.2.
\emph{Manuscript reference:} 57:314-315.
\par Fibrewise Pic0 may not be representable even for smooth projective Enriques families
\par\medskip\noindent\textbf{K.76. SA14-H10-main --- Proved.}\par
Pic0 abelian for smooth projective char0 arbitrary base.
\par\emph{Statement/proof location:} work/bce/additional\_leads.tex; Deligne1968 Thm5.5(i) p123;SGA3 VI\_B Prop1.6 p327.
\emph{Manuscript reference:} 57:64.
\par Deligne Thm5.5 gives conormal=(R1f*O)\textasciicircum{}dual; SGA3 criterion yields smooth Pic0; proper neutral locus yields abelian scheme
\par\medskip\noindent\textbf{K.77. SA14-H10-conormal --- Proved.}\par
Locally free conormal at unit in char0 implies smoothness without flatness.
\par\emph{Statement/proof location:} work/bce/additional\_leads.tex; SGA3 VI\_B Prop1.6 original p327.
\emph{Manuscript reference:} 57:64.
\par For group lfp along identity over Q-base, omega locally free iff smooth there; no flatness assumed
\par\medskip\noindent\textbf{K.78. SA14-H11-flattening --- Proved.}\par
Flattening problem A.
\par\emph{Statement/proof location:} work/bce/additional\_leads.tex; Stacks05PS;0CWQ.
\emph{Manuscript reference:} 57:73-79.
\par Same as E3
\par\medskip\noindent\textbf{K.79. SA14-H11-affinity --- Proved over every integral one-dimensional noetherian base after shrinking.}\par
Replace quasi-affine Pic pullback by affine after shrinking dimension1 base.
\par\emph{Statement/proof location:} work/bce/additional\_leads.tex; SGA3 XVI4.2;VI\_B1.4.2;SGA6 XII1.1-1.2.
\emph{Manuscript reference:} 57:74;78.
\par Full Theorem H11.G via affine flat kernel quotient and fibrewise-finite monomorphism; no normality/projectivity assumption.
\par\medskip\noindent\textbf{K.80. E4-308 --- No mathematical question.}\par
Title leaf.
\par\emph{Statement/proof location:} work/bce/section\_E.tex; transcription.
\emph{Manuscript reference:} 57:308.
\par Title only
\par\medskip\noindent\textbf{K.81. E4-309a --- Refuted.}\par
Pic disjoint sum of finite-type schemes?.
\par\emph{Statement/proof location:} work/bce/section\_E.tex; Kleiman2005 Ex4.14 and deduction.
\emph{Manuscript reference:} 57:309.
\par Semistable split conic has connected component with infinitely many origins
\par\medskip\noindent\textbf{K.82. E4-309aprime --- Refuted.}\par
Closure of quasi-compact subset quasi-compact?.
\par\emph{Statement/proof location:} work/bce/section\_E.tex; Kleiman2005 Ex4.14 and deduction.
\emph{Manuscript reference:} 57:309.
\par Closure of generic point in conic Pic contains infinitely many origins
\par\medskip\noindent\textbf{K.83. E4-309adoubleprime --- Refuted generally.}\par
Pic\ensuremath{\tau} closed in Pic?.
\par\emph{Statement/proof location:} work/bce/section\_E.tex; Kleiman2005 Ex4.14;FGA VI Cor2.3.
\emph{Manuscript reference:} 57:309.
\par Conic Pic tau identity copy is open not closed; normal morphism positive subset result retained separately
\par\medskip\noindent\textbf{K.84. E4-309b --- Refuted in written generality.}\par
Pic/Pic\ensuremath{\tau} quotient scheme?.
\par\emph{Statement/proof location:} work/bce/section\_E.tex; Theorem E4.Q; Kleiman2005 Example4.14 pp31–33.
\emph{Manuscript reference:} 57:309.
\par Nonsplit Mumford conic over R[[t]] has Pic\ensuremath{\tau}=trivial subgroup S but full Pic is a nonschematic algebraic space; quotient equals Pic. Even flat tau is insufficient without closedness. Stronger variants with full Pic assumed schematic or closed nonflat tau are separate.
\par\medskip\noindent\textbf{K.85. E4-309c --- Refuted, including smooth normal families.}\par
Pic\ensuremath{\tau} universally open?.
\par\emph{Statement/proof location:} work/bce/section\_E.tex; section.tex.
\emph{Manuscript reference:} 57:309.
\par E4.H has a clopen character component supported on the closed point. E4.R also refutes openness with reduced fibres by dimension jump.
\par\medskip\noindent\textbf{K.86. E4-309cflat --- Refuted.}\par
Margin: Pic\ensuremath{\tau} flat?.
\par\emph{Statement/proof location:} work/bce/section\_E.tex; PartD.
\emph{Manuscript reference:} 57:309.
\par E4.H gives smooth projective X/W(k), char(k)=2, with Pic\ensuremath{\tau}=\ensuremath{\mu}2\ensuremath{{}^2}×(SpecR disjoint Speck), coordinate module R\textasciicircum{}4\ensuremath{\oplus}k\textasciicircum{}4. Explicit degree768 construction and root audit completed.
\par\medskip\noindent\textbf{K.87. E4-309d --- Proved under standard Pic representability hypotheses.}\par
Pic\ensuremath{\tau} quasi-compact hence finite type?.
\par\emph{Statement/proof location:} work/bce/section\_E.tex; LaurentSchroer2024 Thm2.1.
\emph{Manuscript reference:} 57:309.
\par Open subspace of finite presentation for proper flat fp cohflat0
\par\medskip\noindent\textbf{K.88. E4-309e --- Refuted even when represented.}\par
Pic\ensuremath{\tau} separated when representable?.
\par\emph{Statement/proof location:} work/bce/section\_E.tex; LaurentSchroer2024 Prop2.3(i).
\emph{Manuscript reference:} 57:309.
\par Free K3 involution combined with t=u\ensuremath{{}^2} produces regular projective X with double special fibre; O(Z) has order2 and generic restriction trivial but nontrivial special restriction; Raynaud5.3 ensures full Pic is a scheme
\par\medskip\noindent\textbf{K.89. E4-309bprime --- Mathematical readings settled; pale abbreviation retained.}\par
Semicontinuity of unsure\{Min.|Dim.\} Lissité.
\par\emph{Statement/proof location:} work/bce/section\_E.tex; Official MPL57 scan309; FGA VI Cor2.7 p234; Theorem E4.S; DGJ targeted audit.
\emph{Manuscript reference:} 57:309.
\par Original scan309 prefers Dim. Lissite but remains pale. Fibre dimension is upper semicontinuous; relative smooth locus is open. E4.S gives geometrically integral reduced surface family with nonsmooth generic Pic-tau and smooth special Pic-tau, refuting openness of fibrewise smoothness and upper semicontinuity of tangent-dimension defect. No missing direction is invented.
\par\medskip\noindent\textbf{K.90. E4-309f --- Refuted generally; known positive under geometric normality.}\par
Continuity of ell-rank of Pic fibres.
\par\emph{Statement/proof location:} work/bce/section\_E.tex; FGA VI Cor2.7.
\emph{Manuscript reference:} 57:309.
\par Projective cubic y\ensuremath{{}^2}z=x\ensuremath{{}^3}+x\ensuremath{{}^2}z+tz\ensuremath{{}^3} over C[[t]] has smooth elliptic generic fibre and integral nodal special fibre. Pic[\ensuremath{\ell}] has geometric F\ensuremath{\ell} dimension2 generically and1 specially; normal-family constancy remains FGA VI Cor2.7.
\par\medskip\noindent\textbf{K.91. E4-309g --- Proved under universal H0=O.}\par
Pic\ensuremath{\tau} product formula.
\par\emph{Statement/proof location:} work/bce/section\_E.tex; section.tex forthcoming.
\emph{Manuscript reference:} 57:309.
\par E2 represents Corr by separated unramified scheme; based maps to affine n-torsion prove its geometric fibres torsion-free; fibrewise vanishing descends. Connectedness alone false via dual-number thickening
\par\medskip\noindent\textbf{K.92. E4-309h --- Delegated to root Part A.}\par
Hyperplane Pic\ensuremath{\tau} bijective/isomorphism and surjectivity.
\par\emph{Statement/proof location:} work/bce/section\_E.tex; PartA.
\emph{Manuscript reference:} 57:309.
\par Conditions explicitly left unspecified in manuscript; preserve unsure reading
\par\medskip\noindent\textbf{K.93. E4-3101 --- Proved in represented setting.}\par
Pic\ensuremath{\tau} closed if f normal?.
\par\emph{Statement/proof location:} work/bce/section\_E.tex; FGA VI Cor2.3 p232.
\emph{Manuscript reference:} 57:310.
\par FGA VI Cor2.3; does not imply representability of Pic0
\par\medskip\noindent\textbf{K.94. E4-3102 --- Refuted.}\par
Pic\ensuremath{\tau} flat?.
\par\emph{Statement/proof location:} work/bce/section\_E.tex; PartD.
\emph{Manuscript reference:} 57:310.
\par Same fully specified Heisenberg degeneration as E4.H; independent of the rejected ordinary-Igusa tangent argument.
\par\medskip\noindent\textbf{K.95. E4-3103 --- Refuted in general and for normal families in mixed characteristic.}\par
Pic\ensuremath{\tau} universally open incl normal and char0 cases.
\par\emph{Statement/proof location:} work/bce/section\_E.tex; FGA VI Rem2.9.
\emph{Manuscript reference:} 57:310.
\par E4.H settles the residue-characteristic2 normal-family case negatively. The normal characteristic-zero positive statement remains valid.
\par\medskip\noindent\textbf{K.96. E4-3104 --- Refuted even for a smooth projective family.}\par
Multiplication by every n universally open?.
\par\emph{Statement/proof location:} work/bce/section\_E.tex; section.tex;FGA VI Thm2.5.
\emph{Manuscript reference:} 57:310.
\par E4.H is killed by2, and its identity section is not open: the four geometric generic branches of \ensuremath{\mu}2\ensuremath{{}^2} meet at the special identity. Thus [2] is not open. The cuspidal Ga example gives a simpler field counterexample; invertible n is étale.
\par\medskip\noindent\textbf{K.97. E4-3104ft --- Known with exact theorem.}\par
Margin: multiplication by n finite type?.
\par\emph{Statement/proof location:} work/bce/section\_E.tex; SGA6 XIII Thm3.6 p635;Kleiman Rem6.28.
\emph{Manuscript reference:} 57:310.
\par n nonzero; proper X and represented Pic
\par\medskip\noindent\textbf{K.98. E4-3105ft --- Proved.}\par
Pic\ensuremath{\tau} finite type?.
\par\emph{Statement/proof location:} work/bce/section\_E.tex; LaurentSchroer2024 Thm2.1.
\emph{Manuscript reference:} 57:310.
\par Same as309d
\par\medskip\noindent\textbf{K.99. E4-3105proper --- Refuted.}\par
Pic\ensuremath{\tau} proper for normal morphism?.
\par\emph{Statement/proof location:} work/bce/section\_E.tex; LaurentSchroer2024 Prop2.3(ii).
\emph{Manuscript reference:} 57:310.
\par Explicit family of diagonal K3 three-quadric intersections modulo sign involution: generic Enriques double-cover class of order2 cannot extend across four fixed nodes; all geometric fibres normal integral
\par\medskip\noindent\textbf{K.100. E4-3105projective --- Refuted normal branch; smooth branch positive.}\par
Pic\ensuremath{\tau} projective for smooth or normal projective morphism?.
\par\emph{Statement/proof location:} work/bce/section\_E.tex; Kleiman Ex6.18.
\emph{Manuscript reference:} 57:310.
\par Same normal projective quotient family has nonproper tau, hence nonprojective; smooth projective branch remains the classical positive theorem
\par\medskip\noindent\textbf{K.101. E4-3105prime --- Refuted.}\par
Universal openness implies finite type/proper under normality?.
\par\emph{Statement/proof location:} work/bce/section\_E.tex; transcription.
\emph{Manuscript reference:} 57:310.
\par Normal K3-quotient family in characteristic0 has universally open tau but nonproper tau
\par\medskip\noindent\textbf{K.102. E4-3105diagram --- Refuted for multiplication openness; reading qualifier uncertain.}\par
Margin diagram and multiple-p surjective/open assertion.
\par\emph{Statement/proof location:} work/bce/section\_E.tex; Official MPL57 scan310; enlarged crop MPL57\_310\_diagram.png; Theorem E4.H.
\emph{Manuscript reference:} 57:310.
\par Theorem E4.H explicitly makes multiplication by2 nonopen for a smooth projective family. The mathematical openness assertion is refuted. Original scan310 confirms Non exemple and the branch diagram, but the surrounding surjectivity/component qualifier remains illegible; no additional assertion is invented.
\par\medskip\noindent\textbf{K.103. E4-3106 --- Refuted.}\par
Only finitely many S-points killed by n?.
\par\emph{Statement/proof location:} work/bce/section\_E.tex; section.tex.
\emph{Manuscript reference:} 57:310.
\par Cuspidal cubic Ga over infinite algclosed charp has infinitely many p-torsion points
\par\medskip\noindent\textbf{K.104. E4-3106criterion --- Picard-scope conditions proved; exact iff reading uncertain.}\par
Bracketed criterion for finite type via components over constructible dense opens.
\par\emph{Statement/proof location:} work/bce/section\_E.tex; Official MPL57 scan310; new\_E4\_scope\_310.tex; Laurent-Schroer2.1.
\emph{Manuscript reference:} 57:310.
\par For proper flat fp cohomologically-flat-in-degree-zero families over noetherian S, Pic-tau is fp; both defensible above-open and outside-open finite-component conditions therefore hold for closed irreducible base changes. The exact broader iff quantifiers remain unreadable. An explicit infinite-branch etale group refutes the generic-open component-count criterion for arbitrary lft groups; it is not a Picard realization.
\par\medskip\noindent\textbf{K.105. E4-3107 --- Known with exact theorem.}\par
NS ranks uniformly bounded?.
\par\emph{Statement/proof location:} work/bce/section\_E.tex; SGA6 XIII Thm5.1 p650;Kleiman Rem6.19.
\emph{Manuscript reference:} 57:310.
\par Proper family over noetherian base; rank bounded though not constructible
\par\medskip\noindent\textbf{K.106. E4-3108 --- Refuted.}\par
Connected components of Pic finite type?.
\par\emph{Statement/proof location:} work/bce/section\_E.tex; Kleiman Ex4.14;Cor6.25.
\emph{Manuscript reference:} 57:310.
\par Split conic infinite-origin component; positive for projective geomint family
\par\medskip\noindent\textbf{K.107. E4-3108copies --- Refuted.}\par
Pic disjoint union of copies of Pic\ensuremath{\tau}?.
\par\emph{Statement/proof location:} work/bce/section\_E.tex; Kleiman Ex4.14 and deduction.
\emph{Manuscript reference:} 57:310.
\par Split conic: one connected component has infinitely many origins; tau is one trait
\par\medskip\noindent\textbf{K.108. E4-3109 --- Proved standard quotient.}\par
Definition of Pic/Pic\ensuremath{\tau} over field.
\par\emph{Statement/proof location:} work/bce/section\_E.tex; H6 over field;SGA6 XIII.
\emph{Manuscript reference:} 57:310.
\par Etale group scheme with geometric points NS/torsion; over nonclosed field retain Galois action
\par\medskip\noindent\textbf{K.109. E4-310Howard --- Multiplicity invariance refuted; existence/descent words not a typed statement.}\par
Multiplicity compatibility with specialization; existence/descent.
\par\emph{Statement/proof location:} work/bce/section\_E.tex; PartD;transcription.
\emph{Manuscript reference:} 57:310.
\par E4.H: generic four reduced points versus two special components of length4. Total length4\ensuremath{\to}8, component count4\ensuremath{\to}2, neutral geometric length1\ensuremath{\to}4. Further object for existence/descent is unspecified.
\par\medskip\noindent\textbf{K.110. E4-311open --- Refuted, including smooth normal families.}\par
Pic\ensuremath{\tau} universally open?.
\par\emph{Statement/proof location:} work/bce/section\_E.tex; section.tex.
\emph{Manuscript reference:} 57:311.
\par E4.H directly disproves openness by the special-only character component; E4.R supplies a reduced-fibre dimension-jump example.
\par\medskip\noindent\textbf{K.111. E4-311flat0 --- Refuted in full generality.}\par
Pic0 flat?.
\par\emph{Statement/proof location:} work/bce/section\_E.tex; LaurentSchroer Prop2.2;PartD.
\emph{Manuscript reference:} 57:311.
\par May not even be represented; Part D simple Igusa gives fibre rank jump
\par\medskip\noindent\textbf{K.112. E4-311flatimp --- Refuted.}\par
Pic0 flat implies Pic\ensuremath{\tau} flat?.
\par\emph{Statement/proof location:} work/bce/section\_E.tex; Theorem E4.H;Raynaud1979;explicit character scheme calculation.
\emph{Manuscript reference:} 57:311.
\par Basechange Theorem E4.H to A=W(k)/4. Pic0=\ensuremath{\mu}2\ensuremath{{}^2}\_A flat rank4; Pic\ensuremath{\tau} has additional open component \ensuremath{\mu}2\ensuremath{{}^2}\_k, whose coordinate module k4 is not A-flat. Parent audit passed.; explicit degree768 norm system and first-jet proof, root structural audit accepted
\par\medskip\noindent\textbf{K.113. E4-311canonicalArtin --- Refuted by an explicit dual-number Igusa deformation.}\par
Canonical abelian Pic subgroup on Artin local base?.
\par\emph{Statement/proof location:} work/bce/section\_E.tex; Part D D4a; Katz1981 Theorem2.1(3),(4) pp148–149.
\emph{Manuscript reference:} 57:311.
\par D4a gives k=Fbar2, R=k[epsilon]/epsilon\textasciicircum{}2 and ordinary Serre-Tate parameters q12=1+epsilon, all other qij=1. A compatible abelian Picard-variety subgroup would force 1=q12 by Katz morphism lifting. Pic-tau remains flat. Root, BCE and FHI targeted audits passed; the uncertain manuscript notation is retained.
\par\medskip\noindent\textbf{K.114. E4-311canonicalReduced --- Refuted.}\par
Canonical abelian Pic subgroup on reduced base?.
\par\emph{Statement/proof location:} work/bce/section\_E.tex; Part D D3.
\emph{Manuscript reference:} 57:311.
\par Part D3 cyclic shear gives an explicit reduced trait base with no embedded abelian subgroup having the fibrewise Picard varieties; kernel lengths of (F\textasciicircum{}2,2) change from 2 to 4.
\par\medskip\noindent\textbf{K.115. E4-311NS --- Refuted.}\par
Continuity of torsion of NS.
\par\emph{Statement/proof location:} work/bce/section\_E.tex; Part D D2 and product Igusa computation.
\emph{Manuscript reference:} 57:311.
\par Part D4 product Igusa family has ordinary generic NS torsion Z/2 and supersingular special NS torsion zero; the component-group assertion is false.
\par\medskip\noindent\textbf{K.116. E4-311Igusa --- Proved.}\par
Study modular family of Igusa surface.
\par\emph{Statement/proof location:} work/dgj/section\_D.tex D4o and D4s; section\_D\_supersingular.tex.
\emph{Manuscript reference:} 57:311.
\par Both algebraized miniversal hulls computed: ordinary W[[t,s,u,v]]/(s\ensuremath{{}^2}+2s); supersingular W[[t,s,u,v]]/(s\ensuremath{{}^4}−t*s\ensuremath{{}^2}−2s). Canonical fixed elliptic subgroup and all-Artinian Pic-tau flatness proved; root/fhi independent audits pass.
\par\medskip\noindent\textbf{K.117. E4-312 --- No Picard question.}\par
Tensor-product insertion.
\par\emph{Statement/proof location:} work/bce/section\_E.tex; transcription.
\emph{Manuscript reference:} 57:312.
\par Unrelated Bourbaki tensor-product leaf retained in inventory
\par\medskip\noindent\textbf{K.118. E4-313i --- Proved.}\par
Pic separated if fibres geometrically locally integral?.
\par\emph{Statement/proof location:} work/bce/section\_E.tex; LaurentSchroer Prop2.3(i);Raynaud1970 6.1.9.
\emph{Manuscript reference:} 57:313.
\par Finite-etale Stein reduction makes locally integral connected fibres integral; Pic separation descends
\par\medskip\noindent\textbf{K.119. E4-313ii --- Proved.}\par
Pic\ensuremath{\tau} separated if fibres geometrically reduced?.
\par\emph{Statement/proof location:} work/bce/section\_E.tex; Raynaud1970 Ex6.1.6 p49;Thm6.4.1(3) p56;Cor6.4.5 p57.
\emph{Manuscript reference:} 57:313.
\par Geometrically reduced fibres imply propertyN and dprime=1 on alltraits; Raynaud criterion yields separated tau
\par\medskip\noindent\textbf{K.120. E4-313ntorsion --- Proved under geometrically reduced fibres.}\par
n-torsion Pic separated under same fibre hypothesis?.
\par\emph{Statement/proof location:} work/bce/section\_E.tex; Raynaud1970 Cor6.4.5 and deduction.
\emph{Manuscript reference:} 57:313.
\par n-torsion is a subfunctor of separated tau
\par\medskip\noindent\textbf{K.121. E4-313dim --- Refuted even for geometrically reduced fibres.}\par
Constancy of dim Pic0.
\par\emph{Statement/proof location:} work/bce/section\_E.tex; section\_E.tex Theorem E4.R;SGA3 VI\_B4.2;FGA VI2.7.
\emph{Manuscript reference:} 57:313.
\par Two ruled surfaces glued along a split Poincare ribbon give a projective flat geometrically reduced connected family with universal H0=O and Pic0=E times Ga\textasciicircum{}\{h0(P\_s)\}, dimensions1to2. Normal/no-additive case remains constant.
\par\medskip\noindent\textbf{K.122. E4-313alpha --- Proved in represented proper-flat Picard setting.}\par
Abelian-rank inequality alpha<=alpha prime.
\par\emph{Statement/proof location:} work/bce/section\_E.tex; Artin1969 Thm3.1 p30;GilleGuralnick2023 Thms1.1,4.1;Raynaud1970 3.3.1,5.2.
\emph{Manuscript reference:} 57:313.
\par Unprimed=special. Closure G of generic tau is flat; prime-to-p Kummer torsor lifting makes Pic0\_special/G0\_special unipotent; constant dimension and Gille-Guralnick upper semicontinuity of maximal affine commutative subgroup dimension give alpha\_special<=alpha\_generic.
\par\medskip\noindent\textbf{K.123. E4-313lambda --- Proved in represented proper-flat Picard setting.}\par
Additive-rank inequality lambda>=lambda prime.
\par\emph{Statement/proof location:} work/bce/section\_E.tex; GilleGuralnick2023 Thms1.1,4.1 pp1,7-8;Raynaud1970 3.3.1,4.1,5.2.
\emph{Manuscript reference:} 57:313.
\par Unprimed=special. Gille-Guralnick maximal-unipotent-dimension upper semicontinuity applies directly to separated tau. For represented nonseparated Pic use Raynaud separated quotient and reduced neutral groups.
\par\medskip\noindent\textbf{K.124. E4-313ellrank --- Proved in the stated geometrically reduced setting.}\par
Known inequality 2alpha+mu<=2alpha prime+mu prime.
\par\emph{Statement/proof location:} work/bce/section\_E.tex; FGA VI Remark2.8 original pp234–235, visually verified; Theorem2.5 p233; E4.rank.2.
\emph{Manuscript reference:} 57:313.
\par Unprimed is special and primed generic. Choose ell away from residue characteristic and both NS torsion orders; separated etale Pic[ell] lifts distinct special torsion points to distinct generic ones, giving 2alpha\_special+mu\_special<=2alpha\_generic+mu\_generic. Also valid whenever the relevant prime-to-characteristic torsion subspaces are separated.
\par\medskip\noindent\textbf{K.125. E4-313Corr --- Proved.}\par
Corr exists for geomint proper flat families.
\par\emph{Statement/proof location:} work/bce/section\_E.tex; Murre Thm4 plus E2 proof.
\emph{Manuscript reference:} 57:313.
\par E2 stronger hypotheses on H0 and flatness give separated unramified scheme
\par\medskip\noindent\textbf{K.126. E4-314open --- Proved under standard cohomological representability hypotheses.}\par
Pic\ensuremath{\tau} open of finite presentation.
\par\emph{Statement/proof location:} work/bce/section\_E.tex; LaurentSchroer Thm2.1.
\emph{Manuscript reference:} 57:314.
\par Exact modern theorem; arbitrary Pic functor openness also classical
\par\medskip\noindent\textbf{K.127. E4-314locus --- Proved.}\par
Tau-triviality locus open retrocompact.
\par\emph{Statement/proof location:} work/bce/section\_E.tex; LaurentSchroer Thm2.1.
\emph{Manuscript reference:} 57:314.
\par Pullback of open fp Pic\ensuremath{\tau} by the section representing L
\par\medskip\noindent\textbf{K.128. E4-314closedProjective --- Proved over arbitrary scheme base.}\par
Tau locus open closed for projective geomint family.
\par\emph{Statement/proof location:} work/bce/section\_E.tex; Zheng2009 Lemma3.8 pp313–314; Kleiman2005 Th6.3(h) pp53–54 and B.27 p79; Kollar2022 Lemma11.60 p436; Stacks054F,05F4,055B,04AI.
\emph{Manuscript reference:} 57:314.
\par Theorem E4.T: proper flat fp geometrically integral fibres suffice, without projectivity. Regular Galois alteration, Hodge-index vertical correction, valuation transitivity and negativity prove trait specialization; noetherian approximation gives arbitrary scheme base. Separate root and DGJ AI audits passed.
\par\medskip\noindent\textbf{K.129. E4-314removeProjective --- Proved over arbitrary scheme base.}\par
Remove projectivity from closed tau locus under geomint fibres.
\par\emph{Statement/proof location:} work/bce/section\_E.tex; Zheng2009 Lemma3.8 pp313–314; Kleiman2005 Th6.3(h) pp53–54 and B.27 p79; Kollar2022 Lemma11.60 p436; Stacks054F,05F4,055B,04AI.
\emph{Manuscript reference:} 57:314.
\par Theorem E4.T: proper flat fp geometrically integral fibres suffice, without projectivity. Regular Galois alteration, Hodge-index vertical correction, valuation transitivity and negativity prove trait specialization; noetherian approximation gives arbitrary scheme base. Separate root and DGJ AI audits passed.
\par\medskip\noindent\textbf{K.130. E4-314normal0 --- Proved as subset/valuative statement.}\par
Neutral part closed/proper for normal morphism.
\par\emph{Statement/proof location:} work/bce/section\_E.tex; FGA VI Thm1.1 and Cor2.3.
\emph{Manuscript reference:} 57:314.
\par Not representability of fibrewise neutral subfunctor
\par\medskip\noindent\textbf{K.131. E4-315extension --- Proved intended / refuted literal.}\par
Algebraically trivial generic bundle extends.
\par\emph{Statement/proof location:} work/bce/section\_E.tex; section.tex.
\emph{Manuscript reference:} 57:315.
\par E1 distinction between relative normality and normal total space
\par\medskip\noindent\textbf{K.132. E4-315schemeNormal --- Unsettled arbitrary base.}\par
Pic scheme for proper flat geomnormal fibres.
\par\emph{Statement/proof location:} work/bce/section\_E.tex; Raynaud1970 Thm3.3.1; Theorem E4.Curve; BLR6.6(2c); Cesnavicius2015 provenance.
\emph{Manuscript reference:} 57:315.
\par Projective families, normal noetherian dimension-one bases, and all geometrically normal curve families over arbitrary bases are schematic. The remaining gate is arbitrary-base higher-dimensional geometrically normal families; no genuine Picard counterexample or general proof is obtained. See E4\_schematicity\_attempts.md.
\par\medskip\noindent\textbf{K.133. E4-315schemeIntegral --- Unsettled arbitrary base.}\par
Pic scheme for proper flat geomint fibres.
\par\emph{Statement/proof location:} work/bce/section\_E.tex; Laurent-Schroer2.1; Raynaud1970 Thm3.3.1; Theorem E4.Curve; E4\_schematicity\_attempts.md.
\emph{Manuscript reference:} 57:315.
\par Artin gives a separated algebraic space. Projective families and smooth proper curves over arbitrary bases are schematic. Integral singular curves and higher-dimensional integral families remain. The standard nonsplit conic has reducible geometric special fibre and does not refute this question; arbitrary nonschematic group spaces also do not provide Picard realizations.
\par\medskip\noindent\textbf{K.134. D0-R6 --- Refuted.}\par
The unrestricted full Picard product of R6.
\par\emph{Statement/proof location:} work/dgj/section\_D.tex.
\emph{Manuscript reference:} 57:286.
\par Missing manuscript hypotheses A\_Gamma=0; Poincare line bundle on E x E disproves unrestricted product.
\par\medskip\noindent\textbf{K.135. D0-tau --- Proved under stated hypotheses.}\par
Pic-tau quotient formula for abelian A,B and a constant translation subgroup on B.
\par\emph{Statement/proof location:} work/dgj/section\_D.tex.
\emph{Manuscript reference:} 57:286.
\par All examples satisfy the stated hypotheses; relative fppf sheaf proof includes nonreduced bases.
\par\medskip\noindent\textbf{K.136. D1 --- Proved.}\par
Both signs and variation of quotient Hodge asymmetry.
\par\emph{Statement/proof location:} work/dgj/section\_D.tex.
\emph{Manuscript reference:} 57:287.
\par Characteristic two shearing involutions with F and V; dimensions (2,3),(3,2), specializing to (3,3).
\par\medskip\noindent\textbf{K.137. D2 --- Proved.}\par
A nonflat fixed subgroup and nonflat Pic-tau in a smooth projective quotient family.
\par\emph{Statement/proof location:} work/dgj/section\_D.tex.
\emph{Manuscript reference:} 57:288.
\par Fixed lengths 8 and 16; algebra R[z]/(z\textasciicircum{}4,t z\textasciicircum{}2).
\par\medskip\noindent\textbf{K.138. D3 --- Proved.}\par
No abelian subscheme with fibrewise reduced Picard identity components.
\par\emph{Statement/proof location:} work/dgj/section\_D.tex.
\emph{Manuscript reference:} 57:288.
\par Fibrewise image of (F\textasciicircum{}2,[2]) has kernel degree 2 then 4.
\par\medskip\noindent\textbf{K.139. D4-product --- Proved under stated hypotheses.}\par
Picard scheme and torsion NS in the product Igusa family.
\par\emph{Statement/proof location:} work/dgj/section\_D.tex.
\emph{Manuscript reference:} 57:311.
\par Pic-tau flat; Picvar compatible; NS torsion Z/2 generically and zero at supersingular point.
\par\medskip\noindent\textbf{K.140. D4-versal --- Proved.}\par
Full modular or versal Igusa deformation family.
\par\emph{Statement/proof location:} work/dgj/section\_D.tex D4o and D4s; section\_D\_supersingular.tex.
\emph{Manuscript reference:} 57:311.
\par Both algebraized miniversal hulls computed: ordinary W[[t,s,u,v]]/(s\ensuremath{{}^2}+2s); supersingular W[[t,s,u,v]]/(s\ensuremath{{}^4}−t*s\ensuremath{{}^2}−2s). Canonical fixed elliptic subgroup and all-Artinian Pic-tau flatness proved; root/fhi independent audits pass.
\par\medskip\noindent\textbf{K.141. G2a --- Proved under stated hypotheses.}\par
Determinant boxed relation with explicit sign and typed equivariance.
\par\emph{Statement/proof location:} work/dgj/section\_G.tex.
\emph{Manuscript reference:} 59:10-12.
\par Proved for polarizations, the hypothesis used in Part G.
\par\medskip\noindent\textbf{K.142. G2b --- Proved.}\par
Determinant trivializations agree with canonical sections.
\par\emph{Statement/proof location:} work/dgj/section\_G.tex.
\emph{Manuscript reference:} 59:12-14.
\par Includes square and even-Euler double transform comparisons; corrects ppcm and even-chi divisibility claim.
\par\medskip\noindent\textbf{K.143. G3 --- Proved.}\par
Albanese sign c=-lambda\_Theta and canonical class K=omega\_C.
\par\emph{Statement/proof location:} work/dgj/section\_G.tex.
\emph{Manuscript reference:} 59:23-28.
\par Minus sign visually verified; canonical norm on PPAV theta torsor.
\par\medskip\noindent\textbf{K.144. J0 --- Proved textual scope correction.}\par
Complementary index and distinction of equivalence relations.
\par\emph{Statement/proof location:} work/dgj/section\_J.tex.
\emph{Manuscript reference:} 44:2;44:8.
\par Exponent remains palaeographically uncertain; n+1-i explicit in letter; modern Ab is not silently identified with manuscript quotient. Th1 en substance/ensemblistement and Th4 margin Connu/Annulé or Cancelé transcription conflicts retained.
\par\medskip\noindent\textbf{K.145. J2a --- Refuted.}\par
Universal canonical intersection isomorphism Ab2 to its dual.
\par\emph{Statement/proof location:} work/dgj/section\_J.tex.
\emph{Manuscript reference:} 44:2;44:8-9.
\par Generic monic septic Jacobian over the algebraic closure of F5(a0,...,a6)\textasciicircum{}perf; intersection forces a degree8 K-isogeny, ruled out by irreducible S7 action on A[2]. No integral torsion standard assumption used.
\par\medskip\noindent\textbf{K.146. J2b --- Refuted.}\par
Intersection isomorphism is a polarization up to one sign.
\par\emph{Statement/proof location:} work/dgj/section\_J.tex.
\emph{Manuscript reference:} 44:2;44:8-9.
\par Explicit E/F3 times P1 times P1; matrix [[0,lambda],[lambda,0]].
\par\medskip\noindent\textbf{K.147. F1 --- Known.}\par
Checked prime-to-p, pronilpotent and strongly F-regular finiteness theorems.
\par\emph{Statement/proof location:} work/fhi/F.tex; F\_proof.md.
\emph{Manuscript reference:} 58:19.
\par Full surface group remains undecided
\par\medskip\noindent\textbf{K.148. F2 --- Refuted.}\par
Dimension-free consequence is false for k[[t]] in characteristic p.
\par\emph{Statement/proof location:} work/fhi/F.tex; F\_proof.md.
\emph{Manuscript reference:} 58:19.
\par Correct consequence requires components dimension at least two
\par\medskip\noindent\textbf{K.149. F3 --- Proved.}\par
Metabelian group with pronilpotent quotient Z3 and fg open abelianizations need not be fg.
\par\emph{Statement/proof location:} work/fhi/F.tex; F\_proof.md.
\emph{Manuscript reference:} 58:19.
\par Abstract obstruction; no geometric realization claimed
\par\medskip\noindent\textbf{K.150. F-full --- Undecided.}\par
Full finite generation for all normal complete local surfaces.
\par\emph{Statement/proof location:} work/fhi/F\_proof.md section F4.
\emph{Manuscript reference:} 58:19.
\par For complete normal local surfaces with algebraically closed residue field, the unrestricted conjecture remains. In equal characteristic F8/F9b settle every geometric quotient with normal pro-p kernel over the full tame quotient. F10 proves finite abelian inertia using Krishna–Majumder2026v2; F10c gives all-prime rank-linear H1 bounds for abelian monodromy, and F11 proves uniform finite generation of metabelian quotients. F12 gives the characteristic-p bound for metabelian coefficients; F12b settles pro-p closed second derived subgroup. F13 proves a fixed log-conductor bound after one fixed cover. F13c/d remove unramified character order and reduce the remaining coefficient-prime-not-p estimate to a degree-linear tame-generator bound for cyclic covers of uniformly bounded p-power degree (F13.6). F13e bounds refined-Swan zero length, but not cleanup complexity. F13.5/F13.6, further solvable and nonsolvable wild quotients, and mixed characteristic remain unresolved. The actual wild kernel is not assumed pro-p.
\par\medskip\noindent\textbf{K.151. H1 --- Proved.}\par
Connected pro-generalized Jacobian has projective dimension at most one in Serre category.
\par\emph{Statement/proof location:} work/fhi/H.tex; H\_proof.md.
\emph{Manuscript reference:} 43:24-35;43:80.
\par All proalgebraic targets; all p
\par\medskip\noindent\textbf{K.152. H2 --- Proved.}\par
Ext>=2(P,Ga)=0 for every source P in Serre category.
\par\emph{Statement/proof location:} work/fhi/H.tex; H\_proof.md.
\emph{Manuscript reference:} 43:80.
\par Corrects assertion that Ga remains difficult
\par\medskip\noindent\textbf{K.153. H3 --- Refuted.}\par
Unaugmented all-degree geometric comparison fails in degree zero.
\par\emph{Statement/proof location:} work/fhi/H.tex; H\_proof.md.
\emph{Manuscript reference:} 43:80.
\par Retain augmentation by Z in a specified larger category
\par\medskip\noindent\textbf{K.154. H4 --- Proved.}\par
Etale local cohomology has H3\_m(k[[s,t]],Z/p) nonzero.
\par\emph{Statement/proof location:} work/fhi/H.tex; H\_proof.md.
\emph{Manuscript reference:} 43:70-74;SA19-H4.
\par Refutes topology/category transfer, not original Zariski representability
\par\medskip\noindent\textbf{K.155. I1 --- Proved.}\par
Degree-zero endomorphisms are diagonal in every regular local residual complex.
\par\emph{Statement/proof location:} work/fhi/I.tex; I\_proof.md.
\emph{Manuscript reference:} 43:67.
\par Finite topology gives product adic topology
\par\medskip\noindent\textbf{K.156. I2 --- Proved.}\par
Complete surface presentation including End2=restricted adeles/F.
\par\emph{Statement/proof location:} work/fhi/I.tex; I\_proof.md.
\emph{Manuscript reference:} 43:65-68;SA19-H3.
\par Precise finite-topology universal interpretation; mixed arrow topology
\par\medskip\noindent\textbf{K.157. I-full --- Undecided.}\par
Full presentation under the manuscript completion without a category specification.
\par\emph{Statement/proof location:} work/fhi/I\_universal\_linear.tex; I\_proof.md.
\emph{Manuscript reference:} 43:65-68;SA19-H3.
\par All dimensions are proved for separated complete graded algebras with a basis of open left ideals; all surfaces are proved for arbitrary jointly continuous topology. The original43:66 specifies the broader category without the open-left-ideal hypothesis. I7b now proves open dense differential-image quotients for whole source columns over arbitrary regular local rings. Remaining: construction and continuity of the entire adjacent-arrow family and its endpoint assembly in arbitrary complete targets, in dimension at least3.
\par\medskip\noindent\textbf{K.158. R3g --- Proved.}\par
The generic smooth hyperplane restriction kernel is finite connected.
\par\emph{Statement/proof location:} sections/A.tex, a:generic.
\emph{Manuscript reference:} 57:99-100.
\par Proof uses Kummer Lefschetz, the p-torsion dlog injection, and the incidence variety; targeted independent check by fhi.
\par\medskip\noindent\textbf{K.159. A3f --- Proved.}\par
Direct Raynaud fibre polarizations in the stated degree range are not very ample.
\par\emph{Statement/proof location:} sections/A.tex, a:fibre.
\emph{Manuscript reference:} 57:99-100.
\par Semigroup pole orders force the fibre map to factor through Frobenius; finite check covers 518 pairs, while the proof treats all coprime pairs.
\par\medskip\noindent\textbf{K.160. E4-3103-normal-p --- Refuted.}\par
Pic tau universally open for normal morphism in residue characteristic p.
\par\emph{Statement/proof location:} work/bce/section\_E.tex; Theorem E4.H;Raynaud1979 4.2.1-4.2.6 pp135-139;CotnerZavyalov2024 Thm1.1.8(iii).
\emph{Manuscript reference:} 57:310.
\par Explicit rank8 Heisenberg degeneration over W(k), p2, realized by Raynaud complete intersections; Pic\ensuremath{\tau}=\ensuremath{\mu}2\ensuremath{{}^2} times (SpecR disjoint Speck), so an open component maps only to closed point. Parent audit passed.; explicit degree768 norm system and first-jet proof, root structural audit accepted
\par\medskip\noindent\textbf{K.161. E4-313dim-reduced --- Refuted.}\par
Constancy of Pic dimension with geometrically reduced fibres.
\par\emph{Statement/proof location:} work/bce/section\_E.tex; section\_E.tex Theorem E4.R.
\emph{Manuscript reference:} 57:313.
\par Finite flat double X over a ruled surface, locally z(z-d)=0, glues two reduced components along a doubled section. Universal base-change exactness and units Mayer-Vietoris yield Pic0 dimensions1to2.
\par\medskip\noindent\textbf{K.162. SA14-H11-affinity-normal --- Proved.}\par
Affine Pic pullback after shrinking one-dimensional base, for normal generic fibres.
\par\emph{Statement/proof location:} work/bce/additional\_leads.tex; SGA6 XII1.1;Kleiman4.18.5;Stacks09N9;H6.
\emph{Manuscript reference:} 57:74;78.
\par Full quotient proof in additional\_leads.tex; includes finite/projective/relative-normal candidate and smooth birational candidate
\par\medskip\noindent\textbf{K.163. E4-3105projective-smooth --- Known.}\par
Projectivity of tau for smooth projective morphism.
\par\emph{Statement/proof location:} work/bce/section\_E.tex; LaurentSchroer2.3(ii);Kleiman6.18.
\emph{Manuscript reference:} 57:310.
\par Kleiman2005 Exercise6.18 p59 and Answer6.18 p72 explicitly prove projectivity of tau for smooth projective geometrically integral families over noetherian S; opened primary text. Disconnected smooth fibres reduce via finite étale Stein factor.
\par\medskip\noindent\textbf{K.164. D4-first-order --- Refuted.}\par
Universal first-order Igusa deformation has nonflat Pic-tau.
\par\emph{Statement/proof location:} work/dgj/section\_D.tex D4o; attempt\_D4\_ordinary.tex.
\emph{Manuscript reference:} 57:311.
\par Every Artinian deformation of an Igusa surface, ordinary or supersingular first factor, has flat Pic-tau; full proofs independently audited by root and fhi. Literal Chai2008v2 Theorem2.4(iii) and Chai-Oort2017 first-order assertion are false.
\par\medskip\noindent\textbf{K.165. D4-gap --- Refuted inference; actual Igusa assertion separately refuted by complete proof.}\par
Different absolute tangent dimensions of nonsmooth torsors imply nonflatness.
\par\emph{Statement/proof location:} work/dgj/section\_D.tex D4o; attempt\_D4\_ordinary.tex; finite torsor certificate.
\emph{Manuscript reference:} 57:311.
\par The finite free torsor refutes the tangent inference. Complete ordinary and supersingular deformation calculations additionally refute actual first-order nonflatness; historical unresolved statuses superseded.
\par\medskip\noindent\textbf{K.166. H5 --- Proved.}\par
Higher Ext vanishing also on the perfect etale site with exact category matching.
\par\emph{Statement/proof location:} work/fhi/H.tex; H\_proof.md H5.
\emph{Manuscript reference:} 43:80.
\par Source is extension of abelian variety by affine proalgebraic group; no arbitrary pro-target comparison
\par\medskip\noindent\textbf{K.167. H4b --- Refuted.}\par
Ordinary local Jacobian of k[s,t]\_(s,t) is not projective.
\par\emph{Statement/proof location:} work/fhi/H.tex; H\_proof.md H4b.
\emph{Manuscript reference:} 43:70-74;SA19-H4.
\par Actual nonsplit Z/p extension in Pro(AbSch); all p
\par\medskip\noindent\textbf{K.168. I3 --- Proved.}\par
Universal presentation for every regular local surface via explicit arrow completion.
\par\emph{Statement/proof location:} work/fhi/I.tex; I\_proof.md section I3.
\emph{Manuscript reference:} 43:65-68;SA19-H3.
\par No extra degree-two kernel; arrow completion target identified rigorously
\par\medskip\noindent\textbf{K.169. I4 --- Proved.}\par
Every adjacent residual arrow identifies its full Hom module as a separated coefficient completion.
\par\emph{Statement/proof location:} work/fhi/I.tex; I\_proof.md I4; I\_higher.md.
\emph{Manuscript reference:} 43:65-67;SA19-H3.
\par Topology q\textasciicircum{}N A\_q + p\textasciicircum{}n ÂA\_p; injectivity and density proved; every dimension
\par\medskip\noindent\textbf{K.170. I5 --- Proved.}\par
The diagonal algebra and differential generate all graded endomorphisms densely in every dimension.
\par\emph{Statement/proof location:} work/fhi/I.tex; I\_proof.md I5; I\_higher.md.
\emph{Manuscript reference:} 43:66-67;SA19-H3.
\par Dense generation does not settle quotient topology or all closed relations
\par\medskip\noindent\textbf{K.171. H3b --- Refuted.}\par
Even augmented perfect-site Hom cannot compare naturally to ordinary H0(K,G) for all algebraic groups.
\par\emph{Statement/proof location:} work/fhi/H.tex; H\_proof.md H3b.
\emph{Manuscript reference:} 43:80.
\par Perfection mismatch is independent of missing degree augmentation
\par\medskip\noindent\textbf{K.172. F1b --- Known.}\par
Full punctured etale fundamental group of every complete log terminal surface is finite in all equal and mixed characteristics.
\par\emph{Statement/proof location:} work/fhi/F.tex; F\_proof.md F1b.
\emph{Manuscript reference:} 58:19.
\par Includes p2,p3; log terminality additional; no prime-to-p-order assertion
\par\medskip\noindent\textbf{K.173. A3n --- Proved.}\par
A very ample negative-H1 example requires q\_nil>=h0(L)>=5 on a surface.
\par\emph{Statement/proof location:} sections/A.tex, a:nilpotence.
\emph{Manuscript reference:} 57:99-100.
\par Opened Enokizono Th4.20 with E=0; low projective embeddings have H1(O)=0.
\par\medskip\noindent\textbf{K.174. A3s --- Proved.}\par
An A3 surface is of general type, satisfies explicit numerical bounds, and has no smooth W2 lifting.
\par\emph{Statement/proof location:} sections/A.tex, a:generaltype.
\emph{Manuscript reference:} 57:99-100.
\par Opened Terakawa Th1.6, corrected Zhang Th3.1, and Langer Cor11.2; very ampleness forces degree at least3 on a genus-one fibre.
\par\medskip\noindent\textbf{K.175. A3R --- Proved.}\par
Every positive two-parameter Raynaud polarization with negative H1 fails very ampleness.
\par\emph{Statement/proof location:} sections/A.tex, a:raynaud-all.
\emph{Manuscript reference:} 57:99-100.
\par Connection lemma gives vanishing for a>=ell; fibre semigroup excludes a<ell. This family exclusion is compatible with the separate A3 cover construction.
\par\medskip\noindent\textbf{K.176. C2-DVR --- Proved in all characteristics without flatness of X.}\par
Picard completion isomorphism with finite-length H1 over an actual excellent henselian DVR.
\par\emph{Statement/proof location:} work/bce/section\_C.tex; section\_C.tex Theorem C2.D;Artin1969 Thm3.5.
\emph{Manuscript reference:} 57:272.
\par Complete elementary proof in Theorem C2.D; no H2 finite-length assumption. Flat closure plus Artin-Rees fibre-product patching handles nonflat X.
\par\medskip\noindent\textbf{K.177. SA14-H11-affinity-general --- Proved.}\par
Affine Pic pullback for proper surjections after shrinking a one-dimensional integral base.
\par\emph{Statement/proof location:} work/bce/additional\_leads.tex; SGA3 XVI4.2 pp521-522;VI\_B1.4.2 pp325-326;SGA6 XII1.1-1.2.
\emph{Manuscript reference:} 57:74;78.
\par Kernel quotient and fibrewise-finite monomorphism argument, full proof Theorem H11.G.
\par\medskip\noindent\textbf{K.178. D4-modular-Picvar --- Proved.}\par
A modular ordinary-Igusa degeneration has flat Pic-tau but incompatible fibrewise Picard varieties.
\par\emph{Statement/proof location:} work/dgj/section\_D.tex; section\_D\_modular.tex.
\emph{Manuscript reference:} 57:311; 57:288.
\par Over W(k), E=F:Y2Z+XYZ=X3+XZ2, P=(0:0:1), Q=(-2:1:w) with 2w2-w+8=0. Klein-four free quotient; Pic-tau finite flat degree8 over F-dual; Albanese degrees4 generic and2 special exclude an abelian Picvar subscheme.
\par\medskip\noindent\textbf{K.179. F3b --- Proved.}\par
Explicit metabelian example has H1(G,M\_n)=L\_n\textasciicircum{}n and unbounded normalized twisted cohomology.
\par\emph{Statement/proof location:} work/fhi/F.tex; F\_proof.md F3b.
\emph{Manuscript reference:} 58:19.
\par Necessary d-generator bound dimH1<=d dimM fails; no geometric realization claimed
\par\medskip\noindent\textbf{K.180. I6 --- Proved.}\par
Residual Hom columns are exact after closing images above the source height.
\par\emph{Statement/proof location:} work/fhi/I.tex; I\_proof.md I6; I\_relations.md.
\emph{Manuscript reference:} 43:66-68;SA19-H3.
\par Quadratic kernel identified in actual Hom topology for gap-two intervals; free-algebra topology still separate
\par\medskip\noindent\textbf{K.181. R3u --- Proved.}\par
The generic hyperplane restriction kernel is infinitesimal unipotent for normal projective X.
\par\emph{Statement/proof location:} sections/A.tex, a:generic-unipotent.
\emph{Manuscript reference:} 57:99-100.
\par Rational dlog handles normal X; Raynaud Proposition6.2.3 excludes mu\_p; Cartier duality proves unipotence. Root independently checked the source and proof.
\par\medskip\noindent\textbf{K.182. A-CZ346ii --- Proved.}\par
The additional smooth one-form isomorphism hypothesis in Cotner-Zavyalov Example3.4.6(ii) is incompatible with its alpha\_p kernel.
\par\emph{Statement/proof location:} sections/A.tex, a:cz-vacuity.
\emph{Manuscript reference:} 57:99-100 source route.
\par The conditional auxiliary example is vacuous; this is not a refutation of the paper main theorems and does not settle A1/A2.
\par\medskip\noindent\textbf{K.183. R1 --- Corrected and rechecked.}\par
Supplied divisor counterexample and refined relative-complete-intersection criterion.
\par\emph{Statement/proof location:} sections/B.tex.
\emph{Manuscript reference:} 57:110-111.
\par The displayed affine factorization requires char not2 and i\ensuremath{{}^2}=-1; proper examples in every characteristic supplied. R2 sufficiency retains the manuscript relative lci hypothesis.
\par\medskip\noindent\textbf{K.184. R2 --- Carried as background; not used.}\par
Transgression identity on57:144.
\par\emph{Statement/proof location:} Input readings only.
\emph{Manuscript reference:} 57:144.
\par No result in this run relies on this identity; no new independent verification is claimed.
\par\medskip\noindent\textbf{K.185. R5 --- Carried as background; not used.}\par
Factor2 in the recalled Igusa b2 formula.
\par\emph{Statement/proof location:} Input readings only.
\emph{Manuscript reference:} 57:70.
\par The quotient constructions use directly computed Hodge and fixed-group data, not the supplied b2 formula.
\par\medskip\noindent\textbf{K.186. R6 --- Refuted as stated; corrected in scope.}\par
Unrestricted product formula for the Picard group of a quotient.
\par\emph{Statement/proof location:} sections/D.tex, D0.
\emph{Manuscript reference:} 57:286.
\par Cross-correspondence terms already refute the full Pic product formula for trivial Gamma. The properly scoped Pic-tau formula is proved and used.
\par\medskip\noindent\textbf{K.187. R7 --- Carried as background; not used.}\par
Functional equation determinant identity and correction terms.
\par\emph{Statement/proof location:} Input readings only.
\emph{Manuscript reference:} 60:53.
\par This background claim is not used by PartsA-J proofs in this run; no independent adoption or new verification is claimed.
\par\medskip\noindent\textbf{K.188. R8 --- Corrected and rechecked in stated cases.}\par
Local Lefschetz reductions and finite H1 of open subgroups.
\par\emph{Statement/proof location:} sections/F.tex.
\emph{Manuscript reference:} 58:2-19.
\par Full group finite generation remains open in the run; prime-to-p and pronilpotent results are proved from opened sources. Dimension-one extrapolation is refuted.
\par\medskip\noindent\textbf{K.189. R9 --- Proved with typed signs.}\par
Boxed determinant relation and canonical torsor point criterion.
\par\emph{Statement/proof location:} sections/G.tex, G1 and G2a.
\emph{Manuscript reference:} 59:10-14.
\par Complete criterion and determinant comparison are proved, with quadratic-refinement monodromy obstruction.
\par\medskip\noindent\textbf{K.190. R10 --- Known theorem with scope and date corrections.}\par
Neron component-group pairing and historical comparison.
\par\emph{Statement/proof location:} sections/H.tex, priority appendix.
\emph{Manuscript reference:} 43:117.
\par SGA7IX pairing and its perfectness conjecture are distinguished; current perfect-residue-field theorem cited. Leaf priority is not inferred from dossier dates.
\par\medskip\noindent\textbf{K.191. I7 --- Proved.}\par
Complete regular-ring residual columns are exact and their maps onto cycles are open above the first arrow.
\par\emph{Statement/proof location:} work/fhi/I.tex; I\_complete\_columns.tex; I\_complete\_columns.md.
\emph{Manuscript reference:} 43:66-68;SA19-H3.
\par All dimensions over complete A; uniform universal-family topology remains unresolved
\par\medskip\noindent\textbf{K.192. D4-supersingular --- Proved.}\par
Full supersingular Igusa hull and all-Artinian Pic-tau flatness.
\par\emph{Statement/proof location:} work/dgj/section\_D.tex D4s; section\_D\_supersingular.tex.
\emph{Manuscript reference:} 57:311.
\par Exact ring W[[t,s,u,v]]/(s\ensuremath{{}^4}−t*s\ensuremath{{}^2}−2s), finite free rank4 over W[[t,u,v]]. Pic-tau finite locally free degree8 over F-dual. Complete independent root/fhi audits.
\par\medskip\noindent\textbf{K.193. J1i --- Known, with derived kernel statement Proved.}\par
All-codimension incidence quotient, family kernels, bound and complementary isogeny.
\par\emph{Statement/proof location:} work/dgj/section\_J\_incidence.tex; section\_J.tex.
\emph{Manuscript reference:} 44:2;44:8.
\par Primary1971 record credits Grothendieck1964; all-regular representability not identified. Single-self-pairing shortcut refuted by E×P1\ensuremath{{}^2}.
\par\medskip\noindent\textbf{K.194. I8 --- Proved.}\par
Full residual endomorphism presentation in every dimension for complete algebras with open-left-ideal basis.
\par\emph{Statement/proof location:} work/fhi/I\_universal\_linear.tex.
\emph{Manuscript reference:} 43:64-66.
\par Full proof independently passed by root and DGJ; BCE passed the generator lemma; open-left-ideal category retained; no inference for arbitrary non-linear topology
\par\medskip\noindent\textbf{K.195. E4-309bprime-fibreSmooth --- Refuted.}\par
Is fibrewise Picard smoothness an open condition?.
\par\emph{Statement/proof location:} work/bce/section\_E.tex; new\_E4\_smoothness.tex; code/check\_cuspidal\_igusa.output.json; work/dgj/audit\_E4\_cusp.md.
\emph{Manuscript reference:} 57:309.
\par E4.S: char2 cuspidal Igusa degeneration has Pic-tau=Ydual times G[2], generically Ydual times mu2 times Z2 and specially Ydual times Ga. Original X fibres remain geometrically integral reduced.
\par\medskip\noindent\textbf{K.196. E4-309bprime-defect --- Refuted.}\par
Upper semicontinuity of tangent dimension minus fibre dimension?.
\par\emph{Statement/proof location:} work/bce/section\_E.tex; Theorem E4.S.
\emph{Manuscript reference:} 57:309.
\par Same E4.S has tangent dimension2 on both fibres, geometric dimensions1 generic and2 special; the defect drops1to0.
\par\medskip\noindent\textbf{K.197. E4-309bprime-dimension --- Proved.}\par
Upper semicontinuity of Picard fibre dimension?.
\par\emph{Statement/proof location:} work/bce/section\_E.tex; FGA VI Cor2.7 p234; Stacks fibre dimension theorem.
\emph{Manuscript reference:} 57:309.
\par The relative local dimension function is upper semicontinuous on a locally finitely presented group space; translate to the identity and compose with the identity section.
\par\medskip\noindent\textbf{K.198. E4-309bprime-relativeSmooth --- Proved.}\par
Openness of the relative smooth locus?.
\par\emph{Statement/proof location:} work/bce/section\_E.tex; Standard smooth-locus theorem; E4.S scope distinction.
\emph{Manuscript reference:} 57:309.
\par Smoothness is open for locally finitely presented morphisms. It includes flatness and differs from smoothness of each individual geometric fibre.
\par\medskip\noindent\textbf{K.199. D4-artin-Picvar --- Refuted.}\par
Every Artinian Igusa deformation admits a compatible abelian Picard subscheme.
\par\emph{Statement/proof location:} work/dgj/section\_D.tex D4a; section\_D\_artin\_picvar.tex.
\emph{Manuscript reference:} 57:311.
\par R=Fbar2[epsilon]/epsilon\ensuremath{{}^2}, E0=F0:y\ensuremath{{}^2}+xy=x\ensuremath{{}^3}+1, P=(0,1), q12=1+epsilon, otherq=1. Any abelian Picvar lift would yield f:B\ensuremath{\to}Adual and force1=1+epsilon. Root and bce independently audited; matches FGA VI Remark3.4 p238.
\par\medskip\noindent\textbf{K.200. J6L --- Proved under explicit algebraicity hypotheses.}\par
Conditional Lefschetz theorem for incidence abelian varieties.
\par\emph{Statement/proof location:} work/dgj/section\_J.tex J6L; section\_J\_lefschetz.tex.
\emph{Manuscript reference:} 44:4 conditional stable and high ranges.
\par Cohomologically zero correspondences act trivially on incidence varieties. Lambda\_X j\_* is stable inverse; j\_* Lambda\_Y is high-range right inverse. No extra sl2 coefficients; root independently audited.
\par\medskip\noindent\textbf{K.201. J6P --- Proved under explicit smooth-axis and generic-member hypotheses.}\par
Codimension-two pencil restriction equals the trace image, with stable-range surjectivity.
\par\emph{Statement/proof location:} work/dgj/section\_J\_pencil.tex; section\_J.tex J6P.
\emph{Manuscript reference:} 44:4 middle clause for i=2, n=4; stable surjectivity for i=2, n>=5.
\par All characteristics. Actual Ab2 blowup decomposition absorbs exceptional Picvar B via Picvar X. Finite kernel additionally follows from stated torsion-injectivity assumption, known char0. Independent root audit PASS.
\par\medskip\noindent\textbf{K.202. J6Z --- Proved.}\par
Highest nonzero target codimension restriction Pinc\textasciicircum{}(n-1)X to AlbY is surjective.
\par\emph{Statement/proof location:} work/dgj/section\_J\_pencil.tex; section\_J.tex J6Z.
\emph{Manuscript reference:} 44:4 high range i=n-1 for n>=3.
\par Any smooth hyperplane section in every characteristic. Also universal regular representatives whenever source exists, in particular Ab2 of any smooth projective threefold. Independent root audit PASS.
\par\medskip\noindent\textbf{K.203. F4 --- Proved.}\par
Uniform normalized simple-module H1 bounds imply quantitative finite generation of every prosolvable group.
\par\emph{Statement/proof location:} work/fhi/F\_prosolvable\_criterion.tex; F.tex.
\emph{Manuscript reference:} 58:19.
\par At most ceil(C)+1 generators; geometric uniform bound still unproved; nonsolvable quotients separate
\par\medskip\noindent\textbf{K.204. F1c --- Known.}\par
Mixed-characteristic BCM-regular pairs have finite prime-to-p local fundamental group bounded by inverse perfectoid signature.
\par\emph{Statement/proof location:} work/fhi/F.tex.
\emph{Manuscript reference:} 58:19.
\par Complete normal local ring; algebraically closed residue; effective boundary and Q-Cartier/BCM-regular pair are additional assumptions
\par\medskip\noindent\textbf{K.205. F5 --- Proved.}\par
Uniform rank-linear H1 bound for all abelian prime-to-p monodromy F\_p local systems on complete normal equal-characteristic local surfaces.
\par\emph{Statement/proof location:} work/fhi/F\_tame\_abelian\_bound.tex; F.tex.
\emph{Manuscript reference:} 58:19.
\par Excludes F3 as a geometric quotient in char2; no full prosolvable or full pi1 finite-generation claim
\par\medskip\noindent\textbf{K.206. F6 --- Proved.}\par
Every pro-p-by-abelian pro-prime-to-p quotient of a complete normal equal-characteristic local surface group is uniformly finitely generated.
\par\emph{Statement/proof location:} work/fhi/F\_pro\_p\_by\_abelian.tex; F.tex.
\emph{Manuscript reference:} 58:19.
\par Also virtually abelian tame quotient gives finite generation; arbitrary solvable towers and full group remain undecided
\par\medskip\noindent\textbf{K.207. A2D --- Proved.}\par
For fixed nontrivial p-torsion M and very ample L, smooth hyperplanes trivializing M form a finite determinantal locus, and a nonempty locus requires h1(M)-h1(M(-L))>=h0(L)-2.
\par\emph{Statement/proof location:} sections/A.tex, a:determinantal; work/root/A2\_determinantal.tex.
\emph{Manuscript reference:} 57:97.
\par Complete conormal and maximal-minor proof; BCE independent targeted audit passed. Gives a finite fixed-pair test, not a universal exclusion or a counterexample.
\par\medskip\noindent\textbf{K.208. E4-315schemeCurve --- Proved.}\par
Full Picard scheme for smooth proper geometrically connected curves over arbitrary base.
\par\emph{Statement/proof location:} work/bce/section\_E.tex; Theorem E4.Curve; BLR1990 6.6 Th2(c) pp164–166, 9.4 Prop4 p260; Stacks07S6; prior derivation Cesnavicius2015.
\emph{Manuscript reference:} 57:315.
\par Genus0 or >=2 is projective by the relative canonical bundle. In genus1 C=Pic1 is an E-torsor; Picn=C/E[n] and BLR6.6 Th2(c) gives affine neighborhoods of every finite orbit, so finite-flat quotient is schematic. Noetherian approximation gives arbitrary base.
\par\medskip\noindent\textbf{K.209. F7 --- Proved.}\par
Uniform rank-linear H1 bound for all finite F\_p systems extending over a fixed resolution.
\par\emph{Statement/proof location:} work/fhi/F\_unramified\_coefficients.tex; F.tex.
\emph{Manuscript reference:} 58:19.
\par Includes nonabelian and wild coefficient monodromy; extension over fixed resolution required
\par\medskip\noindent\textbf{K.210. F7a --- Proved.}\par
Relative generation criterion for a normal pro-p kernel over any finitely generated quotient.
\par\emph{Statement/proof location:} work/fhi/F\_unramified\_coefficients.tex; F.tex.
\emph{Manuscript reference:} 58:19.
\par No solvability of quotient required; uniform simple F\_p-module H1 bound required
\par\medskip\noindent\textbf{K.211. F7b --- Proved.}\par
Pro-p extensions of local fundamental-group quotients factoring through a fixed resolution are uniformly finitely generated.
\par\emph{Statement/proof location:} work/fhi/F\_unramified\_coefficients.tex; F.tex.
\emph{Manuscript reference:} 58:19.
\par Quotient may be nonabelian and wild; does not assert actual kernel of pi1(U) to pi1(X) is pro-p
\par\medskip\noindent\textbf{K.212. J6C --- Refuted.}\par
The middle fixed factor equals the image for every integral-fibre hyperplane pencil.
\par\emph{Statement/proof location:} work/dgj/section\_J6\_cubic.tex; section\_J.tex J6C.
\emph{Manuscript reference:} 44:4 middle clause with n=4 and i=2.
\par Characteristic5: X=sum\_\{j=0\}\textasciicircum{}2 x\_j\textasciicircum{}3+x3\textasciicircum{}2*x4+x4\textasciicircum{}3+x5\textasciicircum{}3, pencil x5=t*x4. All fibres geometrically integral. Ab2X=Pinc2X=0. Ruling correspondence from constant Pic0(Fermat cubic) has transpose composite[-2], so target trace is positive. Root and BCE full independent audits PASS; exact finite certificate saved. No priority claim; Eckardt-cone geometry has classical sources.
\par\medskip\noindent\textbf{K.213. D4b --- Refuted.}\par
Every deformation of a Jacobian ordinary Igusa surface lifts its Albanese elliptic fibration.
\par\emph{Statement/proof location:} work/dgj/section\_D\_partsch.tex; section\_D.tex D4b.
\emph{Manuscript reference:} 57:311; associated modern-source verification.
\par Any lift to a smooth proper genus-one curve gives pullback from its abelian Pic0 lifting the special Picvar inclusion, contradicted by q12=1+epsilon. Refutes literal Partsch2013 Proposition6.6 pp102-103, Theorem6.7 p104 and Proposition6.9 p106 for this actual ordinary Jacobian Igusa deformation. Root full audit PASS; finite Z2-to-mu2 map certifies the etale-source inference problem separately.
\par\medskip\noindent\textbf{K.214. JWalker --- Unresolved; exact equivalent dimension test proved.}\par
Walker isogeny-lift universality implies universality for all regular homomorphisms.
\par\emph{Statement/proof location:} work/dgj/section\_J\_walker.tex; section\_J.tex JWalker.
\emph{Manuscript reference:} 44:2 all-regular interpretation of Th1-3 remains a gate.
\par Stronger universality is equivalent to dim image(usual AJ,phi)=dim Ja for every regular test phi. Positive-dimensional kernel of the first projection is not ruled out by divisible AJ kernel. Kahn final Conjecture3.1 p847 and August2026 ACMV version checked; neither conjecture nor general existence settled.
\par\medskip\noindent\textbf{K.215. F8 --- Proved.}\par
Uniform rank-linear H1 bound for all finite F\_p systems tame along a fixed SNC resolution.
\par\emph{Statement/proof location:} work/fhi/F\_tame\_rootstack.tex; F.tex.
\emph{Manuscript reference:} 58:19.
\par Independent root full-source/proof audit and BCE targeted audit passed; full monodromy may have a p-part; no historical priority claim
\par\medskip\noindent\textbf{K.216. F8b --- Proved.}\par
All local quotients with normal pro-p kernel and arbitrary pro-prime-to-p quotient are uniformly finitely generated.
\par\emph{Statement/proof location:} work/fhi/F\_tame\_rootstack.tex; F.tex.
\emph{Manuscript reference:} 58:19.
\par No abelian or solvable hypothesis on prime-to-p quotient; alternating tame/wild normal factors remain outside the conclusion
\par\medskip\noindent\textbf{K.217. F9 --- Known method; proved here.}\par
Uniform generator bound for the full local tame quotient of a fixed SNC resolution.
\par\emph{Statement/proof location:} work/fhi/F\_tame\_quotient.tex; F.tex.
\emph{Manuscript reference:} 58:19.
\par Full root audit passed; d\_X depends only on component genera and exceptional dual graph; no priority claim
\par\medskip\noindent\textbf{K.218. F9b --- Proved.}\par
Local quotients with normal pro-p kernel and quotient factoring through the full tame local group are uniformly finitely generated.
\par\emph{Statement/proof location:} work/fhi/F\_tame\_quotient.tex; F.tex.
\emph{Manuscript reference:} 58:19.
\par Allows tame quotients with p-part and no solvability assumption; actual wild kernel is not assumed pro-p
\par\medskip\noindent\textbf{K.219. A2V --- Proved necessary condition.}\par
A Cartier-fixed regular one-form vanishing along a smooth ample divisor has order pa and a nonzero exact tangential leading term, yielding a Frobenius kernel on H1(Y,O\_Y(-aY)).
\par\emph{Statement/proof location:} sections/A.tex; work/dgj/A2\_cartier\_leading.tex.
\emph{Manuscript reference:} 57:97.
\par Also forces h0(Omega\_X)>=h0(O\_X(pY)) and a nonnormal associated mu\_p torsor. The affine model dlog(1+t\textasciicircum{}p x) disproves naive p-squared order iteration. Root independent audit and finite coefficient checks passed. Full A2 global extension/exclusion remains unresolved.
\par\medskip\noindent\textbf{K.220. A2W --- Proved under explicit closedness hypothesis.}\par
All global one-forms closed implies Picard point injectivity for every smooth very ample divisor in dimension at least3.
\par\emph{Statement/proof location:} work/dgj/A2\_closed\_forms.tex; root integration in Part A.
\emph{Manuscript reference:} 57:97-100 partial exclusion.
\par H0(Omega1 tensor L\textasciicircum{}-1)=0 in dimension at least2. Any point-kernel example in dimension n>=3 must have dim image d:H0Omega1 to H0Omega2 >=n-1; it necessarily has nonclosed global one-forms. The hypothesis is not automatic. Root audited the all-closed proof; quantitative refinement follows by the same wedge rank.
\par\medskip\noindent\textbf{K.221. F10 --- Proved.}\par
Finite kernel of the abelian local fundamental group map to a fixed resolution and uniform rank-one Swan bound.
\par\emph{Statement/proof location:} work/fhi/F\_abelian\_wild\_inertia.tex; F.tex.
\emph{Manuscript reference:} 58:19.
\par Depends on Nov8 2026 preprint in exact regular F-finite category; root full and BCE targeted audits passed; no priority claim
\par\medskip\noindent\textbf{K.222. F10b --- Proved.}\par
One fixed finite abelian cover removes inertia for all original abelian coefficient monodromy systems.
\par\emph{Statement/proof location:} work/fhi/F\_abelian\_wild\_inertia.tex; F.tex.
\emph{Manuscript reference:} 58:19.
\par All primes and ranks; same cover and resolution; no prime-to-coefficient-degree assumption
\par\medskip\noindent\textbf{K.223. F10c --- Proved.}\par
Uniform rank-linear H1 bound over every coefficient prime for abelian monodromy.
\par\emph{Statement/proof location:} work/fhi/F\_abelian\_wild\_inertia.tex; F.tex.
\emph{Manuscript reference:} 58:19.
\par Root full audit passed; fixed-cover restriction need not be injective; finite kernel contribution bounded
\par\medskip\noindent\textbf{K.224. F11 --- Proved.}\par
Every metabelian quotient of the local fundamental group is uniformly finitely generated.
\par\emph{Statement/proof location:} work/fhi/F\_abelian\_wild\_inertia.tex; F.tex.
\emph{Manuscript reference:} 58:19.
\par Derived length two including alternating tame and wild factors; derived length three and arbitrary nonsolvable quotients remain open here
\par\medskip\noindent\textbf{K.225. A2Y --- Proved under either explicit target hypothesis; independent root audit PASS.}\par
Target Frobenius injectivity or a smooth W2 lifting of the target excludes a Picard point kernel.
\par\emph{Statement/proof location:} work/dgj/A2\_target\_frobenius.tex; root integration in Part A.
\emph{Manuscript reference:} 57:97-100 partial exclusion.
\par n>=3, Y smooth very ample. H1(L\textasciicircum{}-a) injects H1(OY) by a smooth connected E in|L\textasciicircum{}a| and multiplication commutes with Frobenius using s\textasciicircum{}p. A W2 lifting of Y alone makes all required negative H1 vanish in total degree1<min(p,dimY); no X/polarization lifting or dimY<p needed.
\par\medskip\noindent\textbf{K.226. F12 --- Proved.}\par
Uniform H1/rank bound for simple coefficient-p systems with metabelian monodromy.
\par\emph{Statement/proof location:} work/fhi/F\_metabelian\_p\_coefficients.tex; F.tex.
\emph{Manuscript reference:} 58:19.
\par Equal characteristic p; root full proof audit passed; restricted system need not be simple and cover degree need not be prime to p
\par\medskip\noindent\textbf{K.227. F12b --- Proved.}\par
All local quotients whose closed second derived subgroup is pro-p are uniformly finitely generated.
\par\emph{Statement/proof location:} work/fhi/F\_metabelian\_p\_coefficients.tex; F.tex.
\emph{Manuscript reference:} 58:19.
\par Second derived subgroup may be nonabelian; full derived-length3 coefficient-ell-not-p gate remains open here
\par\medskip\noindent\textbf{K.228. F13a --- Proved.}\par
Every simple modular representation of an abelian-by-abelian finite group is induced from a character of a subgroup containing the specified abelian normal subgroup.
\par\emph{Statement/proof location:} work/fhi/F\_metabelian\_log\_conductor.tex; F.tex.
\emph{Manuscript reference:} 58:19.
\par DGJ independent proof; root full audit passed; inducing subgroup normal with abelian quotient
\par\medskip\noindent\textbf{K.229. F13b --- Proved.}\par
Rank-one Swan cutoff is invariant under arbitrary connected finite etale covers of the fixed resolution.
\par\emph{Statement/proof location:} work/fhi/F\_metabelian\_log\_conductor.tex; F.tex.
\emph{Manuscript reference:} 58:19.
\par Root and BCE geometry audits passed; complete local Stein normalization checked; no bound on total character order inferred
\par\medskip\noindent\textbf{K.230. F13 --- Proved.}\par
All metabelian coefficient-ell-not-p systems have a fixed highest logarithmic break bound after one fixed cover.
\par\emph{Statement/proof location:} work/fhi/F\_metabelian\_log\_conductor.tex; F.tex.
\emph{Manuscript reference:} 58:19.
\par Root full proof and source audit passed; rank-linear complete-local Betti transfer F13.5 remains open; Hu-Teyssier finite-type theorem not applied beyond its scope
\par\medskip\noindent\textbf{K.231. A2Z --- Proved construction exclusion; independently audited; canonicalized by root integrator.}\par
A finite separable degree-p cover cannot kill a p-torsion class on a smooth ample full inverse image.
\par\emph{Statement/proof location:} work/dgj/A2\_fierce\_cover.tex.
\emph{Manuscript reference:} A2 failed-construction route only.
\par For finite flat generically separable degree p X to B smooth projective, C smooth ample and Y=pi*C smooth integral, Pic(C)[p](k) injects Pic(Y)[p](k). Nonnormal mu\_p torsor case included: effective conductor adds to positive ramification intersection. No closure of general A2.
\par\medskip\noindent\textbf{K.232. A2S-cross --- Known deduction from existing A3s; initial weaker quasi-elliptic formulation superseded.}\par
A threefold point-kernel target is general type and its leading order obeys the A3s numerical bounds.
\par\emph{Statement/proof location:} work/dgj/A2\_target\_surfaces.tex.
\emph{Manuscript reference:} A2 partial target exclusion.
\par For A2V integer a, a\ensuremath{{}^2}L\ensuremath{{}^2}<=B\_p(Y) from existing A3s since L\textasciicircum{}a is very ample. No remaining quasi-elliptic alternative and no new example.
\par\medskip\noindent\textbf{K.233. A2Gs-cross --- Undecided candidate exclusion; independent proof and source audit pending.}\par
No smooth G\_s torsor over a surface with smooth ample full inverse image.
\par\emph{Statement/proof location:} work/dgj/attempt\_A2\_Gs\_smooth\_surface.tex.
\emph{Manuscript reference:} A2 failed-construction route only.
\par No smooth surface Z with smooth ample D and line subbundle O(-D) of TZ tangent to D. Includes G\_s torsors with both total space and full inverse image smooth. Singular normalization and higher dimensions remain outside this argument.
\par\medskip\noindent\textbf{K.234. F13c --- Proved.}\par
Uniform bounded p-power ramified order modulo an unramified twist on every finite etale resolution cover.
\par\emph{Statement/proof location:} work/fhi/F\_bounded\_cyclic\_reduction.tex; F.tex.
\emph{Manuscript reference:} 58:19.
\par Root full proof audit passed; same exponent e on every cover; Swan conductors unchanged by the unramified twist
\par\medskip\noindent\textbf{K.235. F13d --- Proved.}\par
A degree-linear tame-generator bound for cyclic covers of fixed bounded p-power degree suffices for the remaining metabelian coefficient estimate.
\par\emph{Statement/proof location:} work/fhi/F\_bounded\_cyclic\_reduction.tex; F.tex.
\emph{Manuscript reference:} 58:19.
\par Root full proof audit passed; only sufficient reduction proved; geometric bound F13.6 and complete-local Betti bound F13.5 remain Undecided
\par\medskip\noindent\textbf{K.236. F13e --- Proved.}\par
Refined-Swan zero length is bounded linearly in the degree of every finite etale resolution cover.
\par\emph{Statement/proof location:} work/fhi/F\_refined\_swan\_zeros.tex; F.tex.
\emph{Manuscript reference:} 58:19.
\par Root full proof audit passed; component bound d\_a(-B\_ii)(t\_i-s\_a); total at most d sum\_i(-B\_ii)t\_i; no quantitative cleanup or local Betti conclusion inferred
\par\medskip\noindent\textbf{K.237. F13f --- Proved.}\par
A post-cover conductor cutoff below p gives uniform finite generation at derived length three.
\par\emph{Statement/proof location:} work/fhi/F\_small\_conductor\_corollary.tex; F.tex.
\emph{Manuscript reference:} 58:19.
\par Root audit passed; extra hypothesis t\_i<p on the fixed post-F10b resolution; earlier t\_i<1 criterion sharpened; no assertion such a cover always exists
\par\medskip\noindent\textbf{K.238. I7b --- Proved.}\par
Residual Hom-column differentials have open dense image without assuming the regular local ring complete.
\par\emph{Statement/proof location:} work/bce/I\_strict\_dense\_columns.tex; I.tex.
\emph{Manuscript reference:} 43:64-66.
\par BCE proof and root independent audit passed; arbitrary-target extension is conditional on preceding-degree continuity; whole-family continuity and endpoint assembly remain open
\par\medskip\noindent\textbf{K.239. A1Q --- Known; derivation proved.}\par
Picard restriction for quasi-F-split smooth ample targets, and for every smooth Fano threefold target.
\par\emph{Statement/proof location:} sections/A.tex, a:known-targets; work/root/A1\_known\_targets.tex.
\emph{Manuscript reference:} 57:97.
\par Quasi-F-split targets give injectivity for ambient dimension at least3 and isomorphism for dimension at least4. A smooth Fano threefold target gives an isomorphism in every p, including2. Exact target hypotheses and duality indices independently audited by DGJ; no unrestricted A1 or A2 conclusion.
\par\medskip\noindent\textbf{K.240. F13g --- Proved.}\par
A negative Swan intersection forces its positive coefficient to be divisible by p; nonzero Swan divisors have maximum coefficient at least p.
\par\emph{Statement/proof location:} work/fhi/F\_swan\_residue\_gap.tex; F.tex.
\emph{Manuscript reference:} 58:19.
\par Root full proof and exact primary-source audit passed; applies on every finite etale resolution cover
\par\medskip\noindent\textbf{K.241. F13h --- Proved.}\par
The uniform unramified exponent improves to any e with p\textasciicircum{}(e+1) greater than the maximal conductor cutoff.
\par\emph{Statement/proof location:} work/fhi/F\_swan\_residue\_gap.tex; F.tex.
\emph{Manuscript reference:} 58:19.
\par Root audit passed; sharpens the stated sufficient exponent by one p-step where applicable; F13.5 and F13.6 still open
\par\medskip\noindent\textbf{K.242. A2P --- Proved construction exclusion.}\par
A smooth ample relative Frobenius p-multisection of a projective-line bundle cannot kill p-torsion from its base.
\par\emph{Statement/proof location:} sections/A.tex; work/root/A2\_multisection\_exclusion.tex.
\emph{Manuscript reference:} 57:97.
\par For a quotient F*E to N the normal line is Q\textasciicircum{}p and the relative canonical line Q\textasciicircum{}(p-1). Normalization of a mu-p torsor would make the latter anti-effective, contradicting ample normal bundle. Positive-dimensional smooth projective base, quotient convention and full relative-Frobenius inverse section explicit. DGJ independent equation/conductor audit passed. Does not settle A2.
\endgroup
